\documentclass{amsart}
\usepackage{mathtools, amsmath, amssymb, amsthm, url, enumerate}
\usepackage{tikz}
\usetikzlibrary{cd}
%\everymath{\displaystyle}

 
% Page layout
\usepackage[margin=25mm]{geometry}
 
\usepackage[all, cmtip]{xy}
\usepackage[hypertexnames=false]{hyperref}
 
\newcommand{\R}{\mathbb{R}}
\newcommand{\Z}{\mathbb{Z}}
\newcommand{\C}{\mathbb{C}}
\newcommand{\Q}{\mathbb{Q}}
\newcommand{\LH}{\mathbb{H}}
\newcommand{\E}{\mathrm{E}}
\newcommand{\fh}{\mathfrak{h}}
\newcommand{\vol}{\mathrm{vol}}
\newcommand{\Clin}{C^{\infty}_{\mathrm{lin}}}
\newcommand{\X}{\mathcal{X}}
 
\DeclareMathOperator{\Ot}{\mathcal{O}}
\DeclareMathOperator{\D}{\mathcal{D}}
\DeclareMathOperator{\Leg}{Leg}
\DeclareMathOperator{\Iso}{Iso}
\DeclareMathOperator{\Path}{Path}
\DeclareMathOperator{\Img}{Img}
\DeclareMathOperator{\pr}{pr}
\DeclareMathOperator{\ad}{ad}
\DeclareMathOperator{\Diff}{Diff}
\DeclareMathOperator{\Para}{Para}
\DeclareMathOperator{\dom}{dom}
\DeclareMathOperator{\id}{id}
\DeclareMathOperator{\evl}{ev}
\DeclareMathOperator{\Int}{Int}
\DeclareMathOperator{\Map}{Map}
\DeclareMathOperator{\Loop}{Loop}
\DeclareMathOperator{\stPath}{stPath}
\DeclareMathOperator{\Ker}{Ker}
\DeclareMathOperator{\colim}{colim}

 
\theoremstyle{definition}
 
\newtheorem{thm}{Theorem}[section]
\newtheorem{defn}[thm]{Definition}
\newtheorem{rem}[thm]{Remark}
\newtheorem{prop}[thm]{Proposition}
\newtheorem{lem}[thm]{Lemma}
\newtheorem{cor}[thm]{Corollary}
\newtheorem{ex}[thm]{Example}
\newtheorem{prob}[thm]{Open Question}
\newtheorem*{thm*}{Theorem}
\newtheorem*{defn*}{Definition}
\newtheorem{conj}[thm]{Conjecture}
\newtheorem{notation}[thm]{Notation}
\newtheorem{convention}[thm]{Convention}
 
\title{On Riemannian Geometry in Elastic Diffeology}
\author[Y. Shiobara]{Yusuke Shiobara}
\address{%
  Department of Science and Technology,
  Graduate School of Medicine, Science and Technology,
  Shinshu University,
  Matsumoto, Nagano 390-8621, Japan
}
\email{24hs603d@shinshu-u.ac.jp}

\subjclass[2020]{Primary 58A40; Secondary 18F15, 18F40, 53C05, 53C22, 58D15, 58B20}
\keywords{diffeology, metric, tangent functor, connection, mapping space.}

\begin{document}
\begin{abstract}
  Several inequivalent notions of tangent space exist for diffeological spaces, which has hindered the development of Riemannian geometry in this setting. 
  On Blohmann's elastic diffeological spaces, the tangent functor carries a genuine tangent structure in the sense of Cockett and Cruttwell, and we use this to develop Riemannian geometry.
  We introduce connections adapted to this setting and show that they induce covariant derivatives on vector fields satisfying the usual axioms, together with the associated curvature tensor, whose basic properties we record following Cockett--Cruttwell. Given a Riemannian metric, we establish a Koszul formula and prove that the Levi-Civita covariant derivative is unique if it exists, while its existence remains open in general.
  We then develop the theory of geodesics and of the energy functional: geodesics are always critical paths of the energy functional, and on elastic spaces for which the real line is a curve object,  completeness of the geodesic spray yields existence and uniqueness of geodesics, as well as the converse statement that critical paths are geodesics. 
  As a fundamental class of examples, we show that for a closed manifold $M$ and an elastic Riemannian diffeological space $N$ these structures on the mapping space $C^\infty(M,N)$ are
  obtained pointwise from those on $N$, under a natural tangent-commuting hypothesis that holds automatically when $N$ is a manifold; in particular, the geodesic spray of $C^\infty(M,N)$ is complete whenever that of $N$ is.
\end{abstract}
\maketitle

% ============================================================
\section{Introduction}\label{sec:intro}
% ============================================================

Diffeology, introduced by Souriau and systematically developed by Iglesias-Zemmour~\cite{PIZ12}, is a natural generalization of smooth manifolds. The category of diffeological spaces contains that of smooth manifolds as a full subcategory while enjoying closure properties that the latter lacks: it is complete and cocomplete, and mapping spaces are again diffeological spaces. It therefore provides a broad and flexible framework that naturally accommodates singular spaces and infinite-dimensional manifolds.

On this foundation, the topological side of the theory is by now well developed: differential forms and de~Rham theory, sheaf-theoretic methods, and homotopy theory~\cite{Kihara} all admit satisfactory extensions to diffeological spaces. Differential geometry, by contrast, has not yet reached a comparable stage, and the same is true of Riemannian geometry. One important reason is that several inequivalent notions of tangent space have been proposed for diffeological spaces, and these do not generally agree~\cite{Taho}. As a consequence, geometric constructions that depend on a tangent structure lack a canonical starting point.

Two recent developments have substantially advanced the study of Riemannian geometry on diffeological spaces. Iglesias-Zemmour~\cite{PIZ25} proposed a notion of Riemannian metric that avoids tangent spaces altogether, thereby circumventing the ambiguity arising from competing tangent-space constructions. Subsequently, Kuribayashi, Sakai, and the author~\cite{KSS} showed that this metric theory admits a natural tangent-space interpretation through Blohmann's tangent functor~\cite{B23}, constructed as a colimit of the tangent bundles of domains. This provides a tangent-space interpretation of the metric of Iglesias-Zemmour and places it within a broader categorical framework.

While these developments provide satisfactory notions of Riemannian metric, they do not yet furnish the full apparatus of Riemannian geometry. Connections, curvature, geodesics, and variational principles require not merely a tangent functor but a tangent structure in the sense of tangent category theory, compatible with addition, scalar multiplication, and the canonical flip. Blohmann's elastic diffeological spaces~\cite{B23} were introduced precisely so that his tangent functor carries such a structure.

Riemannian geometry on diffeological spaces has also been studied by Goldammer and Welker~\cite{GW}. Motivated by applications to optimization, they introduced a Riemannian metric and a Levi-Civita connection on Vincent's tangent spaces~\cite{V}, a construction distinct from Blohmann's tangent functor. Their Levi-Civita connection is formulated directly as a covariant derivative on vector fields, whereas our approach starts from a connection map in the sense of tangent category theory and derives the corresponding covariant derivative. Nevertheless, the covariant derivative induced by our connection satisfies the same fundamental axioms as in the classical theory (see Propositions~\ref{prop:torsion-free-bracket}, \ref{prop:covariant-derivative-properties}, and \ref{prop:metric-compatibility-vf}). Thus, one major difference between the two frameworks lies in their underlying tangent-space constructions. A detailed comparison is left for future work. We note, however, that in their setting as well, only the uniqueness of the Levi-Civita connection is established, whereas its existence remains open.

The purpose of this paper is to develop Riemannian geometry on elastic diffeological spaces---Levi-Civita connections, curvature, geodesics, and the variational theory of the energy functional---within the framework of the tangent categories of Cockett and Cruttwell~\cite{CC14,CC17}.

More precisely, we introduce a notion of connection on elastic diffeological spaces following Cockett--Cruttwell~\cite{CC17}, supplemented by two scalar axioms reflecting the $\R$-multiplication available in elastic diffeology (Remark~\ref{rem:scalar-axiom}). We show that such a connection induces a covariant derivative on vector fields enjoying the familiar properties from classical differential geometry: $C^\infty(M)$-linearity in the first argument, the Leibniz rule, and, in the Levi-Civita case, metric compatibility and torsion-freeness.

Given a Riemannian metric, we establish a Koszul formula (Proposition~\ref{prop:koszul-formula}) and prove that the Levi-Civita covariant derivative is unique if it exists (Theorem~\ref{thm:uniqueness-cov-deriv}). In contrast with the classical situation, the existence of a Levi-Civita connection is not automatic (see Open Question~\ref{prob:open-questions}).


We further develop the geodesic and variational theory. Here a genuinely infinite-dimensional phenomenon arises: $\mathsf{Elast}$ contains the category of convenient vector spaces, in which smooth dynamical systems may fail to have unique solutions, so that $\R$ is not a curve object in the sense of Cockett--Cruttwell--Lemay~\cite{CCL} in $\mathsf{Elast}$. To isolate the setting in which uniqueness is recovered, we introduce the full subcategory $\mathsf{Elast}_{\mathrm{curve}}$ of elastic spaces for which $\R$ is a curve object (Definition~\ref{def:elast-lc}). Membership amounts to the uniqueness of solutions of parametrized dynamical systems---the categorical remnant of the Picard--Lindel\"of theorem---and every finite-dimensional or Banach smooth manifold belongs to this subcategory.

On $\mathsf{Elast}_{\mathrm{curve}}$, completeness of the geodesic spray yields existence and uniqueness of geodesics (Theorem~\ref{thm:geodesic-existence}). The associated geodesic flow plays the role of an exponential map and realizes variation fields (Lemma~\ref{lem:realization}).

We then establish the variational theory. Geodesics are always critical paths of the energy functional (Proposition~\ref{prop:geodesic-critical-path}), and on $\mathsf{Elast}_{\mathrm{curve}}$, for a full connection whose vertical part is Levi-Civita and whose geodesic spray is complete, the converse also holds (Proposition~\ref{prop:critical-geodesic}). We furthermore prove the first variation formula (Lemma~\ref{lem:first-variation-formula}), from which it follows that length-minimizing paths are geodesics after constant-speed reparametrization (Proposition~\ref{prop:minimizer-geodesic}).

The curve-object hypothesis is robust in our main class of examples. In particular, for mapping spaces we obtain existence and uniqueness of geodesics on $C^\infty(M,N)$ whenever $N \in \mathsf{Elast}_{\mathrm{curve}}$, without deciding whether $C^\infty(M,N)$ itself belongs to $\mathsf{Elast}_{\mathrm{curve}}$ (Remark~\ref{rem:mapping-space-lc-membership}).

Finally, mapping spaces provide a large and natural class of examples. For a closed manifold $M$ and an elastic Riemannian diffeological space $N$, we show that the tangent structure, connections, Riemannian metrics, Levi-Civita connections, curvature, and geodesics on $C^\infty(M,N)$ are obtained pointwise from those on $N$ under a natural tangent-commuting hypothesis, automatic when $N$ is a manifold. In particular, the geodesic spray of $C^\infty(M,N)$ is complete whenever that of $N$ is (Theorem~\ref{thm:summary-of-mapping-space}).

\subsection*{Organization}
The paper is organized as follows. Section~\ref{sec:tangent-structure} reviews Blohmann's tangent functor and the elastic axioms and establishes the basic differential-geometric tools used throughout the paper. Section~\ref{sec:connections} develops the theory of connections and covariant derivatives on elastic diffeological spaces. Section~\ref{sec:riemannian-metric} introduces Riemannian metrics and Levi-Civita structures and proves the corresponding uniqueness results. Section~\ref{sec:curvature} recalls the curvature of a vertical connection from \cite{CC17} and records its basic properties in the elastic setting, with a view toward future applications. Section~\ref{sec:geodesics} develops the geodesic and variational theory. Section~\ref{sec:mapping-space} treats mapping spaces as a fundamental class of examples. Section~\ref{sec:perspectives} discusses directions for future research. Appendix~\ref{sec:tangent-categories} collects the necessary background on tangent categories.

\subsection*{Setting and conventions}
Throughout the paper, $\mathsf{Elast}$ denotes Blohmann's category of elastic diffeological spaces~\cite{B23} and $T$ his tangent functor, which equips $\mathsf{Elast}$ with the structure of a Cartesian tangent category (Section~\ref{sec:tangent-structure}). 
Smooth manifolds are regarded as elastic diffeological spaces via their standard diffeology, and $T$ then restricts to the classical tangent functor. Composition is written in the applicative order, $g \circ f$; note that \cite{CC17, CCL} use the diagrammatic order. Connections are understood in the sense of Cockett--Cruttwell~\cite{CC17}, supplemented by the scalar axioms of Remark~\ref{rem:scalar-axiom}, which reflect the $\R$-module structure of Blohmann's tangent spaces. 
Riemannian metrics are definite by definition (Definition~\ref{def:metric}); where only a weak metric in the sense of \cite{KSS} is needed, we say so explicitly. Standing assumptions specific to a section or subsection --- notably membership in $\mathsf{Elast}_{\mathrm{curve}}$ for the geodesic theory of Section~\ref{sec:geodesics}, and the tangent-commuting hypothesis for mapping spaces in Section~\ref{sec:mapping-space} --- are stated at the head of the relevant part.


\subsection*{Acknowledgement}
I would like to express my sincere gratitude to my supervisor, Keiichi Sakai, for his valuable guidance and continuous support since my undergraduate years. I am also grateful to my sub-advisor, Katsuhiko Kuribayashi, for many insightful discussions and helpful advice. I would like to thank David Miyamoto for his generous support and encouragement.

This work was supported by JST SPRING, Grant Number JPMJSP2144, at Shinshu University.

\subsection*{Declaration of generative AI in the writing process}
During the preparation of this manuscript, the author used ChatGPT and Claude for language editing and for discussing possible formulations, presentations, and future research directions. All mathematical results, proofs, and scientific conclusions were developed and verified by the author. The author carefully reviewed and revised all AI-generated suggestions and takes full responsibility for the content of this publication.


 
% ============================================================
\section{Tangent Structure on Diffeological Spaces}\label{sec:tangent-structure}
% ============================================================

\subsection{Elastic Diffeological Spaces}

Throughout this paper, $\mathsf{Diff}$ denotes the category of diffeological spaces and
smooth maps; we assume familiarity with plots, inductions, subductions, and the
D-topology \cite{PIZ12}. The purpose of this subsection is to fix notation for
Blohmann's diffeological tangent functor and for the class of \emph{elastic}
spaces on which it carries a tangent structure in the sense of
Rosick\'y--Cockett--Cruttwell. All results in this subsection are due to Blohmann
\cite{B23,B24}; we state them without proof, in the form and notation in which
they will be used in Sections~\ref{sec:connections}--\ref{sec:mapping-space}.
A detailed exposition, to which we are indebted, can also be found in
\cite[Section~2.1]{M}.

We denote by $\mathsf{Euc}$ the category whose objects are open subsets of
Euclidean spaces and whose morphisms are smooth maps, and let
$\mathcal{Y}\colon \mathsf{Euc} \to \mathsf{Diff}$ be the full and faithful
functor assigning to each domain its standard diffeology. On $\mathsf{Euc}$ one
has the classical tangent functor
$\hat{T}U = U \times \R^{\dim U}$,
$\hat{T}f(u,v) = (f(u), D_uf(v))$.

\begin{defn}[{\cite[Section~4.1]{B23}}]\label{def:tangent-functor}
For an endofunctor $\hat{F}\colon \mathsf{Euc} \to \mathsf{Euc}$ and a natural
transformation $\hat{\alpha}\colon \hat{F} \to \hat{G}$, set
\begin{equation*}
  \mathbb{L}\hat{F} \coloneqq \mathrm{Lan}_{\mathcal{Y}}\,\mathcal{Y}\hat{F},
  \qquad
  \mathbb{L}\hat{\alpha} \coloneqq
  \mathrm{Lan}_{\mathcal{Y}}\,\mathcal{Y}\hat{\alpha}.
\end{equation*}
The \textbf{tangent functor} $T \colon \mathsf{Diff} \to \mathsf{Diff}$ is the
left Kan extension $T \coloneqq \mathbb{L}\hat{T}$.
\end{defn}

Because $\mathsf{Euc}$ is small and $\mathsf{Diff}$ is cocomplete, these Kan
extensions are computed pointwise: regarding the diffeology $\D(M)$ of a space
$M$ as a category (objects: plots $P\colon U_P \to M$; morphisms: smooth maps
$h\colon U_P \to U_Q$ with $Q \circ h = P$), one has
\begin{equation*}
  (\mathbb{L}\hat{F})M \;=\; \colim_{P \in \D(M)} \hat{F}U_P ,
  \qquad\text{in particular}\qquad
  TM \;=\; \colim_{P \in \D(M)} \bigl(U_P \times \R^{\dim U_P}\bigr).
\end{equation*}

\begin{rem}\label{rem:tangent-structure-2}
We denote by $[P, w]$ the class in $(\mathbb{L}\hat{F})M$ of an element
$w \in \hat{F}U_P$; thus a tangent vector is written $[P,u,v]$ with
$(u,v) \in U_P \times \R^{\dim U_P}$. The identification is the equivalence
relation \emph{generated} by
\begin{equation*}
  [P \circ h,\; w] \;=\; [P,\; \hat{F}h(w)]
  \qquad\text{for every morphism of plots } h ,
\end{equation*}
so that $[P,w] = [P',w']$ precisely when the two representatives are joined by
a finite chain of such elementary moves (in either direction) over
intermediate plots of $M$. On classes, the Kan extension of a natural
transformation acts representative-wise,
$(\mathbb{L}\hat{\alpha})_M[P,w] = [P, \hat{\alpha}_{U_P}(w)]$, and the tangent
map of a smooth map $f \colon M \to N$ is
\begin{equation*}
  T(f)(\,[P,u,v]) = [\,f \circ P,\, u,\, v\,].
\end{equation*}
\end{rem}

The Euclidean tangent bundle carries its usual structure maps
$\hat{\pi}(u,v) = u$, $\hat{0}(u) = (u,0)$, and
$\hat{\kappa}(t,(u,v)) = (u,tv)$, all natural in $U$; their Kan extensions
therefore exist on all of $\mathsf{Diff}$. For $\hat{\kappa}$ a word of caution
is in order, since its source $\R \times \hat{T}$ involves a product: the
identification $\bigl(\mathbb{L}(\R \times \hat{T})\bigr)M \cong \R \times TM$
used below is legitimate because $\mathsf{Diff}$ is Cartesian closed, so that
$\R \times (-)$ is a left adjoint and preserves the colimit defining $TM$. By
contrast, the \emph{fibre} products entering the addition and the flip do not
commute with the colimit; this is the source of the difficulties addressed by
the elasticity axioms below.

\begin{defn}\label{def:nat-trans}
The \textbf{projection} $\pi \coloneqq \mathbb{L}\hat{\pi}$,
\textbf{zero section} $0 \coloneqq \mathbb{L}\hat{0}$, and
\textbf{scalar multiplication} $\kappa \coloneqq \mathbb{L}\hat{\kappa}$ act
on classes by
\begin{equation*}
  \pi_M[P,u,v] = P(u), \qquad
  0_M(x) = [P,u,0] \quad (\text{any } P \text{ with } P(u)=x), \qquad
  \kappa_M(t,[P,u,v]) = [P,u,tv].
\end{equation*}
The value $0_M(x)$ is independent of the chosen plot: if $P(u) = P'(u') = x$,
then both classes agree with $[\,\mathrm{const}_x,\,0,\,0\,]$ by naturality of
$\hat{0}$ applied to the constant plots factoring through $P$ and $P'$.
\end{defn}

The tangent bundle is generated by curves in the following sense; this is the
reason why, in the variational theory of Section~\ref{sec:geodesics}, tangent
vectors to path spaces can be realized by genuine variations.

\begin{defn}\label{def:delta}
  Define $\delta = \delta_M \colon C^{\infty}(\mathcal{Y}\R, M) \to TM$ by
  $\delta(\gamma) \coloneqq [\gamma, 0, 1]$.
\end{defn}

\begin{prop}\label{prop:tangent-structure-1}
  The functor $T$ shares the following properties with the tangent functor on
  manifolds:
  \begin{itemize}
    \item $T$ preserves finite products: the natural map
      $T(M \times N) \to TM \times TN$ is an isomorphism
      {\cite[Corollary~2.2.19]{B24}};
    \item $T$ preserves subductions {\cite[Proposition~2.2.16]{B24}};
    \item $T$ is local: for $U \subseteq M$ D-open, $T$ of the inclusion is an
      induction with D-open image, so $TU$ is an open subspace of $TM$
      {\cite[Proposition~3.35]{B23}};
    \item $\delta_M$ is a subduction {\cite[Proposition~2.2.20]{B24}}.
  \end{itemize}
\end{prop}

Two further pieces of structure are needed for a tangent category: the
fibrewise addition, whose Euclidean model $\hat{+}$ is defined on the fibre
product $\hat{T}_2 = \hat{T} \times_{\hat\pi} \hat{T}$, and the canonical flip
$\hat{\tau}(u,v_0,v_1,v_{01}) = (u,v_1,v_0,v_{01})$ on the second tangent
functor $\hat{T}^2$, together with the vertical lift
$\hat{\lambda}(u,v) = (u,0,0,v)$. As announced above, the Kan extension does
\emph{not} interact well with these constructions: writing
$T_k \coloneqq T \times_\pi \cdots \times_\pi T$ for the $k$-fold fibre
product, the spaces $\mathbb{L}\hat{T}_k M$ and $T_k M$, and likewise
$\mathbb{L}(\hat{T}^2)M$ and $T^2M$, differ in general, so
$\mathbb{L}\hat{+}$ and $\mathbb{L}\hat{\tau}$ do not immediately descend to
the intended domains. What always exists are the comparison morphisms
\begin{alignat*}{2}
  \theta_k&\colon \mathbb{L}\hat{T}_k \to T_k , \quad
    & \theta_{k,M}[P, u, v_0, \ldots, v_{k-1}]
      &= \bigl([P, u, v_0], \ldots, [P, u, v_{k-1}]\bigr), \\
  \theta^2&\colon \mathbb{L}\hat{T}^2 \to T^2 , \quad
    & \theta^2_M[P, u, v_0, v_1, v_{01}] &= T^2P\,(u, v_0, v_1, v_{01}), \\
  \nu_k&\colon TT_k \to T^2 \times_T \cdots \times_T T^2 , \quad
    & \nu_{k,M}\bigl[(P_1,\ldots,P_k), w\bigr]
      &= \bigl([P_1, w], \ldots, [P_k, w]\bigr).
\end{alignat*}

\begin{ex}\label{ex:non-elastic}
The comparison maps genuinely fail to be isomorphisms. Let
$M = \R/(\Z/2)$ be the quotient of the real line by the reflection
$t \mapsto -t$, equipped with the quotient diffeology. Every tangent vector of
$M$ at the singular point is represented by a plot factoring through the
quotient map, and the two ``branches'' cannot be separated after passing to the
colimit; as a consequence $\theta_{2,M}$ is not injective, so the fibrewise
addition $\mathbb{L}\hat{+}$ does not descend to $T_2M$
(see \cite[Example~3.11]{M} for the details). This is precisely the phenomenon
that the elasticity axioms exclude.
\end{ex}

Blohmann's elasticity axioms single out exactly the spaces on which these
comparison maps behave as they do for manifolds.

\begin{defn}[{\cite[Definition~4.1]{B23}}]\label{def:tangent-structure-1}
  A diffeological space $M$ is \textbf{elastic} if:
  \begin{itemize}
    \item[(E1)] $\theta_{k,M}$ is an isomorphism for every $k \geq 1$;
    \item[(E2)] there is a morphism $\tau_M \colon T^2M \to T^2M$ (the
      \textbf{canonical flip}) with
      $\tau_M \circ \theta^2_M = \theta^2_M \circ (\mathbb{L}\hat{\tau})_M$;
    \item[(E3)] the \textbf{vertical lift}
      $\lambda_M \coloneqq \theta^2_M \circ (\mathbb{L}\hat{\lambda})_M
      \colon TM \to T^2M$ is an induction;
    \item[(E4)] $\nu_{k,M}$ is injective for every $k \geq 1$;
    \item[(E5)] each iterated space $T_{k_1}\cdots T_{k_n}M$ satisfies
      (E1)--(E4).
  \end{itemize}
  Elastic spaces span a full subcategory
  $\mathsf{Elast} \subseteq \mathsf{Diff}$.
\end{defn}

\begin{thm}[{\cite[Theorem~4.2]{B23}}]\label{thm:tangent-structure-2}
  The functor $T$, the ring $\R$, and the natural transformations $\pi$, $0$,
  $\kappa$, $+ \coloneqq \mathbb{L}\hat{+} \circ \theta_2^{-1}$, the flip
  $\tau$ of \textup{(E2)}, and the lift $\lambda$ of \textup{(E3)} make
  $\mathsf{Elast}$ a Cartesian tangent category with scalar
  $\R$-multiplication.
\end{thm}

\begin{rem}\label{rem:negatives}
In a tangent category the fibres of $\pi_M$ are a priori only commutative
monoids. Here the scalar $\R$-multiplication supplies negatives: for
$v \in TM$ one has
$v + \kappa_M(-1, v) = 0_M(\pi_M(v))$, as one checks on representatives using
\eqref{eq:TP-linear} below. Hence every $\pi_M \colon TM \to M$, and likewise
every $\pi_{TM} \colon T^2M \to TM$, is a bundle of abelian \emph{groups}, and
differences of morphisms with the same base point are defined. We use this
silently from Section~\ref{sec:connections} on.
\end{rem}

The class of elastic spaces is sufficiently broad for our purposes: it contains
all finite-dimensional manifolds, all convenient vector spaces---in particular
all Banach and Fr\'echet spaces---and the manifolds modelled on them
\cite[Section~2.1]{M}, and it is stable under the mapping-space constructions
of Section~\ref{sec:mapping-space} \cite{B23}. Let us emphasize at once that
this breadth has a price. Smooth dynamical systems on a convenient vector space
need not have unique solutions, so $\mathsf{Elast}$ is strictly larger than the
subcategory $\mathsf{Elast}_{\mathrm{curve}}$ on which the geodesic theory of
Section~\ref{sec:geodesics} takes place; the two should be kept carefully
apart. From now on, unless stated otherwise, every diffeological space is
assumed elastic, and we freely use the structure maps of
Theorem~\ref{thm:tangent-structure-2}.

\begin{rem}\label{rem:cocone-TP}
Since $\mathcal{Y}$ is full and faithful, the left Kan extension along
$\mathcal{Y}$ restricts on representables to the original functor; in
particular $T(\mathcal{Y}U) \cong \mathcal{Y}(\hat{T}U)$ for every
$U \in \mathsf{Euc}$. Under this identification the cocone map of the colimit
$TM = \colim_{P \in \D(M)} \hat{T}U_P$,
\begin{equation*}
  \hat{T}U_P = U_P \times \R^{\dim U_P} \longrightarrow TM,
  \qquad (u,v) \longmapsto [P,u,v],
\end{equation*}
is precisely the tangent map $T(P)$ of the plot $P$: applying the formula of
Remark~\ref{rem:tangent-structure-2} to the identity plot of $U_P$ gives
$T(P)\,[\id_{U_P},u,v] = [P,u,v]$. We henceforth write $T(P)$ for the cocone
maps.
\end{rem}

\begin{lem}\label{lem:TP-linear}
Let $M$ be elastic. Each cocone map $T(P) \colon \hat{T}U_P \to TM$ is a
linear bundle morphism (Definition~\ref{def:differential-bundle-morphism})
over $P$; explicitly, on representatives,
\begin{equation}\label{eq:TP-linear}
  \kappa_M\bigl(r,[P,u,v]\bigr) = [P,u,rv],
  \qquad
  [P,u,v_0] +_M [P,u,v_1] = [P,u,v_0+v_1].
\end{equation}
Moreover the family $\{T(P)\}_{P \in \D(M)}$ is jointly epic, and so is the
family $\{T^{(k)}(P)\}$ of cocone maps of
$\mathbb{L}\hat{T}_kM = \colim_P \hat{T}_kU_P$.
\end{lem}

\begin{proof}
Compatibility with $\pi$ and $0$ is immediate from
Definition~\ref{def:nat-trans}. The first identity in \eqref{eq:TP-linear} is
the naturality of $\kappa$ applied to $P \colon U_P \to M$, and the second is
the naturality of $+$, which is available because $M$ is elastic and hence
$\theta_{2,M}$ is invertible by (E1). Joint epimorphy holds because in
$\mathsf{Diff}$ a colimit is a quotient of a coproduct, so the cocone maps are
jointly surjective and the colimit diffeology is the pushforward of the
diffeologies of the $\hat{T}U_P$.
\end{proof}


% ----------------------------------------
\subsection{Vector Fields, One-Forms and the Lie Bracket}
% ----------------------------------------

This subsection follows the general framework developed in
\cite[Sections~2.4, 2.6]{AB} and \cite[Section~A.1]{M}, where the statements
are established for an arbitrary Cartesian tangent category with scalar
$\R$-multiplication and therefore apply verbatim to $\mathsf{Elast}$ by
Theorem~\ref{thm:tangent-structure-2}. In contrast with
Section~\ref{sec:tangent-structure}.1, the results below that are specific to
the diffeological setting are proved in full, since they are used repeatedly in
the sequel.

\begin{defn}[Vector field]\label{def:vector-field}
Let $M \in \mathsf{Elast}$.
A \textbf{vector field} $V \in \mathfrak{X}(M)$ on $M$ is a smooth section
$V \colon M \to T(M)$, that is, a smooth map satisfying $\pi_M \circ V = \id_M$.
\end{defn}

\begin{defn}\label{def:vf-module}
For $V, W \in \mathfrak{X}(M)$ and $f \in C^\infty(M,\R)$ we set
\begin{equation*}
  V + W \coloneqq {+}_M \circ \langle V, W\rangle,
  \qquad
  fV \coloneqq \kappa_M \circ \langle f, V \rangle ,
\end{equation*}
where $\langle -,- \rangle$ denotes the induced morphism into the fibre
product, respectively into $\R \times TM$. Both are again smooth sections,
since $+_M$ and $\kappa_M$ are morphisms of $\mathsf{Elast}$ over $M$.
\end{defn}

\begin{lem}\label{lem:vf-module}
With the operations of Definition~\ref{def:vf-module}, $\mathfrak{X}(M)$ is a
module over the commutative $\R$-algebra $C^\infty(M,\R)$; in particular it is
an $\R$-module, the scalar $r \in \R$ acting through the constant function $r$.
\end{lem}

\begin{proof}
All axioms are equalities of morphisms $M \to TM$, and by
Lemma~\ref{lem:TP-linear} they may be verified on representatives $[P,u,v]$,
where they reduce to the corresponding identities in the fibres
$\R^{\dim U_P}$ of $\hat{T}U_P$. Negatives are supplied by
Remark~\ref{rem:negatives}.
\end{proof}

We shall use the $C^\infty(M,\R)$-module structure of
Lemma~\ref{lem:vf-module} throughout Sections~\ref{sec:connections}
and~\ref{sec:riemannian-metric}, where the tensoriality of connections and
metrics is expressed in terms of it.

Before defining the action of vector fields on functions we identify the
fibrewise-linear functions on $TM$ with the $1$-forms of Iglesias-Zemmour.

\begin{lem}\label{lem:cotangent-omega}
Let $M$ be an elastic diffeological space, and write $\Clin(-,\R)$ for the set
of fibrewise-linear smooth maps. Then the assignment
$\phi\mapsto(\phi\circ T(P))_P$ is a bijection
\[
  \Clin(TM,\R)
  \;\xrightarrow{\ \sim\ }\;
  \lim_{P\in\D(M)}\Clin(\hat{T}U_P,\R)
  =\lim_{P\in\D(M)}\Omega^{1}(U_P)
  =\Omega^{1}(M),
\]
the limits being taken in the category of sets and the last equality being the
definition of a differential $1$-form on a diffeological space \cite[Section~1.1]{B23}.
\end{lem}

\begin{proof}
Since $C^{\infty}(-,\R)$ is contravariant and carries the colimit
$TM=\colim_P\hat{T}U_P$ to a limit,
\begin{equation}\label{eq:nonlinear-iso}
  C^{\infty}(TM,\R)\;\xrightarrow{\ \sim\ }\;
  \lim_{P}C^{\infty}(\hat{T}U_P,\R),
  \qquad
  \phi\longmapsto(\phi\circ T(P))_{P},
\end{equation}
is a bijection, so it suffices to show that $\phi$ is fibrewise linear if and
only if $\phi_P:=\phi\circ T(P)$ is fibrewise linear for every $P$.

If $\phi$ is fibrewise linear, then so is each $\phi_P$, since $T(P)$ is a
linear bundle morphism (Lemma~\ref{lem:TP-linear}).

Conversely, assume every $\phi_P$ is fibrewise linear. For additivity we must
show $\phi\circ+_M=+_{\R}\circ(\phi\times_M\phi)$ as maps $T_2M\to\R$. Since
$+_M=\mathbb{L}\hat{+}\circ\theta_{2,M}^{-1}$ with
$\theta_{2,M}\colon\mathbb{L}\hat{T}_2M\xrightarrow{\sim}T_2M$ the elastic
isomorphism~(E1), this is equivalent to
$\phi\circ\mathbb{L}\hat{+}=+_{\R}\circ(\phi\times_M\phi)\circ\theta_{2,M}$
on $\mathbb{L}\hat{T}_2M=\colim_P\hat{T}_2U_P$. The cocone maps of this colimit
are jointly epic (Lemma~\ref{lem:TP-linear}), so it suffices to check the
equality on representatives:
\[
  \phi\circ\mathbb{L}\hat{+}\,[P,u,v_0,v_1]
  =\phi_P(u,v_0+v_1)
  =\phi_P(u,v_0)+\phi_P(u,v_1)
  =+_{\R}\circ(\phi\times_M\phi)\,\theta_{2,M}[P,u,v_0,v_1].
\]
Homogeneity is verified in the same manner, on representatives of $TM$
directly: by \eqref{eq:TP-linear} and homogeneity of $\phi_P$,
\[
  \phi\bigl(\kappa_M(r,[P,u,v])\bigr)
  =\phi_P(u,rv)
  =r\,\phi_P(u,v)
  =r\,\phi[P,u,v].
\]
Thus \eqref{eq:nonlinear-iso} restricts to the asserted bijection.
\end{proof}

The categorical definition of the action of a vector field on a function uses
only the tangent structure; Proposition~\ref{prop:pr2-Tf} below identifies it
with the classical formula.

\begin{defn}\label{def:vf-action}
For a function $f \colon M \to \R$ and a vector field $V \in \mathfrak{X}(M)$,
the \textbf{action} of $V$ on $f$ is the function $V \cdot f \colon M \to \R$
defined by
\begin{align*}
    V \cdot f := \pr_2 \circ T(f) \circ V ,
\end{align*}
where $\pr_2 \colon T\R = \R \times \R \to \R$ is the second projection.
\end{defn}

\begin{prop}\label{prop:pr2-Tf}
  Let $d \colon \Omega^n(M) \to \Omega^{n+1}(M)$ denote the exterior derivative.
  For any smooth function $f \colon M \to \R$ one has
  $\pr_2 \circ T(f) = d(f)$
  under the identification of Lemma~\ref{lem:cotangent-omega}; consequently
  $V \cdot f = d(f) \circ V$ for every $V \in \mathfrak{X}(M)$.
\end{prop}

\begin{proof}
  By definition, $d(f) \in \Omega^1(M) = \lim_P \Omega^1(U_P)$ is the family
  $(d(f \circ P))_{P \in \D(M)}$, and for a function $g$ on a domain one has
  $d(g) = \pr_2 \circ \hat{T}g$. Applying
  Remark~\ref{rem:tangent-structure-2} and $T\R = \hat{T}\R$,
  \begin{equation*}
    \bigl(\pr_2 \circ T(f)\bigr)\circ T(P)\,(u,v)
      = \pr_2\bigl[f \circ P, u, v\bigr]
      = \pr_2 \circ \hat{T}(f \circ P)(u,v)
      = d(f \circ P)_u(v)
  \end{equation*}
  for every plot $P$. Thus the image of $\pr_2 \circ T(f) \in \Clin(TM,\R)$
  under the bijection of Lemma~\ref{lem:cotangent-omega} is the family
  $(d(f\circ P))_P = d(f)$, and injectivity of that bijection gives
  $\pr_2 \circ T(f) = d(f)$.
\end{proof}

The action of vector fields on functions has the following properties.

\begin{prop}[{\cite[Proposition~2.25]{AB}}]\label{prop:vf-action-properties}
  For any $V, W \in \mathfrak{X}(M)$ and $f, g \in C^\infty(M, \R)$,
  the following hold:
  \begin{enumerate}
    \item $V \cdot (f + g) = V \cdot f + V \cdot g$;
    \item $(V + W) \cdot f = V \cdot f + W \cdot f$;
    \item $V \cdot (fg) = (V \cdot f)g + f(V \cdot g)$.
  \end{enumerate}
\end{prop}

We now turn to the Lie bracket, for which we first record the standard
extension of the vertical lift.

\begin{defn}\label{def:vertical-lift-extension}
Using the addition of the tangent structure as a bundle of abelian groups
(Remark~\ref{rem:negatives}), the \textbf{extended vertical lift} is the
morphism $\lambda_2 \colon T_2 \to T^2$ defined by
\begin{align*}
    \lambda_2 := \tau \circ {+}_T \circ (T0 \times_0 \lambda),
\end{align*}
where $T0 \times_0 \lambda$ denotes the map induced on the fibre product
$T_2 = T\times_\pi T$ by $T0$ on the first factor and $\lambda$ on the second;
component-wise,
$\lambda_{2,M} = \tau_M \circ {+}_{TM} \circ (T0_M \times_{0_M} \lambda_M)$, so
that informally $\lambda_{2,M}(v,w) = \tau_M(T0_M(v)) +_{TM} \lambda_M(w)$.
\end{defn}

\begin{prop}[{\cite[Section~3.4]{CC14}}]\label{prop:lambda2-iso}
For every $M \in \mathsf{Elast}$ the morphism $\lambda_{2,M}$ is an
isomorphism from $T_2(M) = T(M) \times_M T(M)$ onto the vertical tangent bundle
$\Ker(T\pi_M) \subseteq T^2M$.
\end{prop}

\begin{proof}
This is the universality of the vertical lift, one of the axioms of a tangent
category (Definition~\ref{def:tangent-structure-axioms} in
Appendix~\ref{sec:tangent-categories}), combined with
Theorem~\ref{thm:tangent-structure-2}.
\end{proof}

As in the case of manifolds, one can now define a Lie bracket on vector fields
\cite[Section~3.4]{CC14}.
Let $V, W$ be vector fields on $M$.
By naturality of $\pi_{TM}$,
\begin{equation*}
  \pi_{TM} \circ T(W) \circ V = W \circ \pi_M \circ V = W \circ \id_M = W.
\end{equation*}
Moreover, using the tangent category axiom $\pi_{TM} \circ \tau_M = T(\pi_M)$,
\begin{equation*}
  \pi_{TM} \circ \tau_M \circ T(V) \circ W
    = T(\pi_M) \circ T(V) \circ W
    = T(\id_M) \circ W
    = \id_{TM} \circ W
    = W.
\end{equation*}
Since both expressions have the same base point, their difference is defined in
the bundle of abelian groups $\pi_{TM}\colon T^2M \to TM$
(Remark~\ref{rem:negatives}). This yields the map
$\ell(V,W) \colon M \to T^2(M)$:
\begin{equation*}
  \ell(V,W) = T(W) \circ V - \tau_M \circ T(V) \circ W .
\end{equation*}
Applying the chain rule once more, one finds
\begin{equation*}
  T(\pi_M) \circ \ell(V,W) = V - V = 0_M,
\end{equation*}
so that $\ell(V,W)$ takes values in $\Ker(T\pi_M)$, the vertical tangent
bundle. By Proposition~\ref{prop:lambda2-iso} this kernel is exactly the image
of $\lambda_{2,M}$. Since $\pi_{TM} \circ \ell(V,W) = W$ by the computations
above, the morphism $\ell(V,W)$ factors uniquely through $\lambda_{2,M}$ with
first component $W$; its second component is the vector field we now name.

\begin{defn}\label{def:lie-bracket}
The \textbf{Lie bracket} of vector fields $V, W \in \mathfrak{X}(M)$
is defined as the unique vector field $[V,W] \in \mathfrak{X}(M)$ satisfying
\begin{equation*}
  \ell(V,W) = \lambda_{2,M} \circ \langle W, [V,W] \rangle,
\end{equation*}
where $\langle W, [V,W] \rangle \colon M \to T(M) \times_M T(M)$
is the induced morphism into the fibre product.
\end{defn}

Rosick\'{y} \cite{R} announced that the bracket defined above satisfies
the Jacobi identity; Cockett and Cruttwell \cite{CC15} provided a complete
proof.

\begin{prop}[{\cite[Theorem~4.2]{CC15}}]\label{thm:jacobi}
In any tangent category $\mathcal{C}$, the bracket of
Definition~\ref{def:lie-bracket} is $\R$-bilinear and antisymmetric, and it
satisfies the Jacobi identity for every object $M$.
\end{prop}

Consequently, $\mathfrak{X}(M)$ is an $\R$-Lie algebra. The bracket is
moreover compatible with the action on functions and with the
$C^\infty(M,\R)$-module structure of Lemma~\ref{lem:vf-module}.

\begin{prop}[{\cite[Propositions~2.25, 2.27]{AB}}]\label{prop:leibniz-bracket}
For any $f \colon M \to \R$ and $V, W \in \mathfrak{X}(M)$,
\begin{align*}
    [V, W] \cdot f = V \cdot (W \cdot f) - W \cdot (V \cdot f),
    \qquad
    [V, fW] = (V \cdot f)W + f[V, W].
\end{align*}
\end{prop}

% ============================================================
\section{Connections and Covariant Derivatives}\label{sec:connections}
% ============================================================

Throughout this section, $M$ denotes an elastic diffeological space, and we
freely use the tangent structure of Theorem~\ref{thm:tangent-structure-2}.

% ----------------------------------------
\subsection{Vertical (Affine) Connections and Covariant Derivatives}%
\label{subsec:vertical-cov}
% ----------------------------------------

Connections in tangent categories are described by morphisms on the second tangent bundle rather than directly by covariant derivatives. This viewpoint is particularly suitable in the diffeological setting, where tangent spaces need not admit local trivializations and many classical constructions are unavailable. We begin with the notion of a vertical connection and then derive the corresponding covariant derivative.


\begin{defn}[{\cite[Definitions~3.2, 3.3]{CC17}, modified}]\label{def:vertical-connection}
  A morphism $K \colon T^2M \to TM$ satisfying the following conditions
  is called a \textbf{vertical (affine) connection} on $M$:
  \begin{enumerate}[(a)]
    \item $K \circ \lambda_M = \id_{TM}$
    \item $\pi_M \circ K = \pi_M \circ \pi_{TM}$
    \item $\lambda_M \circ K = T(K) \circ \lambda_{TM}$
    \item $\lambda_M \circ K = T(K) \circ \tau_{TM} \circ T(\lambda_M)$
    \item $K \circ \kappa_{TM} = \kappa_M \circ (\id_{\R} \times K)$
    \item $K \circ T_{(2)}\kappa_M = \kappa_M \circ (\id_{\R} \times K)$
  \end{enumerate}
\end{defn}

\begin{rem}\label{rem:K-linear-bundle}
  By definition, the vertical connection $K$ yields linear bundle morphisms
  (Definition~\ref{def:differential-bundle-morphism}):
  \[
    (K, \pi_M) \colon (T(\pi_M) \colon T^2M \to TM) \to (\pi_M \colon TM \to M),
  \]
  \[
    (K, \pi_M) \colon (\pi_{TM} \colon T^2M \to TM) \to (\pi_M \colon TM \to M).
  \]
  Consequently, $K$ is additive.
\end{rem}

\begin{rem}\label{rem:scalar-axiom}
  In general tangent structures as considered in \cite{CC14, CC15, CC17},
  scalar multiplication on bundles is not part of the axioms.
  In elastic diffeology, however, the structure of scalar $\R$-multiplication
  is present. Taking into account compatibility with connections on tangent spaces
  of manifolds, we have added preservation of scalar multiplication to the axioms
  for both vertical and horizontal connections: axioms~(e) and~(f) for vertical connections,
  and axioms~(e) and~(f) of Definition~\ref{def:horizontal-connection} for
  horizontal connections below. These axioms ensure that the induced covariant derivative
  is linear with respect to scalar multiplication and satisfies the classical Leibniz identities.
\end{rem}

\begin{rem}\label{rem:K-subduction}
  Since $K$ is a split epimorphism, it is a strong epimorphism.
  Strong epimorphisms in the category $\mathsf{Diff}$ of diffeological spaces
  are subductions, so $K$ is a subduction.
\end{rem}

\begin{ex}[Vertical connections on Euclidean space]\label{ex:euclidean-vertical-connection}
Let $M = \R^n$, identified with $TM = \R^n \times \R^n$ and
$T^2M = \R^n \times \R^n \times \R^n \times \R^n$, with coordinates
$(u, v_0, v_1, v_{01}) \in T^2\R^n$ so that
$\pi_{TM}(u, v_0, v_1, v_{01}) = (u, v_0)$ and
$T(\pi_M)(u, v_0, v_1, v_{01}) = (u, v_1)$.
For a smooth family $\Gamma(u)$ of bilinear maps
$\R^n \times \R^n \to \R^n$, the morphism
\begin{align*}
    K_\Gamma(u, v_0, v_1, v_{01}) = \bigl(u,\; v_{01} + \Gamma(u)(v_0, v_1)\bigr)
\end{align*}
is a vertical connection (see \cite[Example~3.6]{CC17}). The bilinearity of
$\Gamma(u)$ is essential: the scalar axioms (e) and (f) amount to the
homogeneity of $\Gamma(u)$ in each of its two arguments.
The special case $\Gamma = 0$, namely
$K_0(u, v_0, v_1, v_{01}) = (u, v_{01})$, is called the
\emph{canonical affine vertical connection}.
Moreover, since $\tau_M(u, v_0, v_1, v_{01}) = (u, v_1, v_0, v_{01})$, we have
$K_\Gamma \circ \tau_M = K_\Gamma$ if and only if each $\Gamma(u)$ is
symmetric; thus torsion-freeness (Definition~\ref{def:torsion-free} below)
recovers the classical symmetry of the Christoffel symbols.
\end{ex}


\begin{defn}[{\cite[Definition~3.15]{CC17}}]\label{def:covariant-derivative}
  Let $K$ be a vertical connection on $M$ and let $X, Y \in \mathfrak{X}(M)$.
  The operation
  \begin{align*}
    \nabla^K \colon \mathfrak{X}(M) \times \mathfrak{X}(M) \to \mathfrak{X}(M),\qquad
    \nabla^K_X Y := K \circ T(Y) \circ X
  \end{align*}
  is called the \textbf{covariant derivative for $K$}.
\end{defn}

\begin{rem}\label{rem:cov-deriv-vs-connection}
  In ordinary Riemannian geometry, connections and covariant derivatives
  determine each other. In a tangent category this need not be the case:
  in the absence of local triviality, the family of morphisms
  $T(Y) \circ X$, for varying $X, Y \in \mathfrak{X}(M)$, need not be
  jointly epic on $T^2M$, so the covariant derivative $\nabla^K$ need not
  determine the vertical connection $K$.
\end{rem}

\begin{defn}[{\cite[Definition~3.24]{CC17}}]\label{def:torsion-free}
  A vertical connection $K$ on $M$ is said to be \textbf{torsion-free} if
  \begin{align*}
    K \circ \tau_M = K.
  \end{align*}
\end{defn}

\begin{prop}[{\cite[Corollary~3.32]{CC17}}]\label{prop:torsion-free-bracket}
  If a vertical connection $K$ on $M$ is torsion-free, then for all
  $X, Y \in \mathfrak{X}(M)$,
  \begin{align*}
    \nabla^K_X Y - \nabla^K_Y X = [X, Y].
  \end{align*}
\end{prop}

\begin{prop}\label{prop:covariant-derivative-properties}
  Let $K$ be a vertical connection on $M$,
  let $X, Y, Z \in \mathfrak{X}(M)$, and let $f \colon M \to \R$ be a smooth
  function. The covariant derivative $\nabla^K$ satisfies the following:
  \begin{enumerate}
    \item $\nabla^K_{X+Y}Z = \nabla^K_X Z + \nabla^K_Y Z$
    \item $\nabla^K_{fX}Y = f\nabla^K_X Y$
    \item $\nabla^K_X (Y + Z) = \nabla^K_X Y + \nabla^K_X Z$
    \item $\nabla^K_X (fY) = (Xf)Y + f\nabla^K_X Y$.
  \end{enumerate}
\end{prop}

\begin{proof}
  \begin{enumerate}
    \item Since $K$ and $T(Z)$ are additive,
    \begin{align*}
      \nabla^K_{X+Y} Z
        = K \circ T(Z) \circ (X + Y)
        = K \circ T(Z) \circ X + K \circ T(Z) \circ Y
        = \nabla^K_X Z + \nabla^K_Y Z.
    \end{align*}
    \item By naturality of the scalar multiplication $\kappa$ and
      the (additional) axiom (e) for vertical connections,
    \begin{align*}
      \nabla^K_{fX}Y
        &= K \circ T(Y) \circ \kappa_M(f, X)
         = K \circ \kappa_{TM}(f, T(Y)\circ X) \\
        &= \kappa_M(f, K \circ T(Y) \circ X)
         = f\nabla^K_X Y.
    \end{align*}
    \item Since the tangent functor preserves addition and $K$ is additive,
    \begin{align*}
      \nabla^K_X (Y + Z)
        = K \circ T(Y+Z) \circ X
        = K \circ T(Y) \circ X + K \circ T(Z) \circ X
        = \nabla^K_X Y + \nabla^K_X Z.
    \end{align*}
    \item By the definition of the covariant derivative,
      $\nabla^K_X (fY) = K \circ T(\kappa_M \circ (f, Y)) \circ X$.
      By functoriality of $T$ and
      Proposition~\ref{prop:decomposition-tangent-morphism},
    \begin{align*}
     & K \circ T(\kappa_M \circ (f, Y)) \circ X \\
        &= K \circ \bigl(T_{(1)} \kappa_M \circ (\id_{T\R} \times \pi_{TM})
            \circ \chi_{\R,TM}
           + T_{(2)} \kappa_M \circ (\pi_{\R} \times \id_{TM})
            \circ \chi_{\R,TM}\bigr)
          \circ (T(f)\circ X, T(Y) \circ X) \\
        &= K \circ T_{(1)} \kappa_M \circ (T(f)\circ X, \pi_{TM} \circ T(Y) \circ X)
           + K \circ T_{(2)} \kappa_M \circ (\pi_{\R} \circ T(f)\circ X, T(Y) \circ X) \\
        &= K \circ T_{(1)} \kappa_M \circ (T(f)\circ X, Y \circ \pi_M \circ X)
           + K \circ T_{(2)} \kappa_M \circ (f \circ \pi_M \circ X, T(Y) \circ X) \\
        &= K \circ T_{(1)} \kappa_M \circ (T(f)\circ X, Y)
           + K \circ T_{(2)} \kappa_M \circ (f, T(Y) \circ X).
    \end{align*}
    We examine each term.
    For $K \circ T_{(1)} \kappa_M \circ (T(f)\circ X, Y)$,
    the second diagram in Definition~\ref{def:scalar-multiplication} gives
    \begin{align*}
      & K \circ T_{(1)} \kappa_M \circ (T(f)\circ X, Y) \\
        &= K \circ \lambda_{2, M} \circ
          \bigl(\kappa_M \circ (\pr_1, \pr_3),\; \kappa_M \circ (\pr_2, \pr_3)\bigr)
          \circ (T(f)\circ X, Y) \\
        &= K \circ \tau_M \circ {+}_{TM} \circ (T0_M \times_{0_M} \lambda_M)
          \circ \bigl(\kappa_M \circ (\pr_1 \circ T(f) \circ X, Y),\;
                 \kappa_M \circ (\pr_2 \circ T(f) \circ X, Y)\bigr) \\
        &= K \circ \tau_M \circ \bigl(0 + \lambda_M \circ \kappa_M
           \circ (\pr_2 \circ T(f) \circ X, Y)\bigr) \\
        &= K \circ \tau_M \circ \lambda_M \circ \kappa_M \circ (df \circ X, Y),
    \end{align*}
    where the last equality uses $\pr_2 \circ T(f) = d(f)$
    (Proposition~\ref{prop:pr2-Tf}).
    Using $\tau_M \circ \lambda_M = \lambda_M$
    (Section~\ref{sec:RosickysAxioms}) and axiom (a),
    we obtain
    \begin{align*}
      K \circ T_{(1)} \kappa_M \circ (T(f)\circ X, Y)
        = \kappa_M \circ (df \circ X, Y) = (Xf)Y,
    \end{align*}
    the last equality by Definition~\ref{def:vf-action}.
    For $K \circ T_{(2)} \kappa_M \circ (f, T(Y) \circ X)$,
    the axiom (f) in Definition~\ref{def:vertical-connection} gives
    \begin{align*}
     K \circ T_{(2)} \kappa_M \circ (f, T(Y) \circ X)
       &= \kappa_{M} \circ  (\id_\R \times K) \circ (f, T(Y) \circ X)\\
       &= \kappa_{M} \circ (f, K \circ T(Y) \circ X) =  f\nabla^K_X Y.
    \end{align*}
    Therefore $\nabla^K_X (fY) = (Xf)Y + f\nabla^K_X Y$.
  \end{enumerate}
\end{proof}

\begin{ex}\label{ex:covariant-derivative}
 Let $M = \R^n$ with the vertical connection $K_\Gamma$ of
Example~\ref{ex:euclidean-vertical-connection}. Write vector fields as
$X(x) = (x, \bar{X}(x))$ and $Y(x) = (x, \bar{Y}(x))$ with
$\bar{X}, \bar{Y} \colon \R^n \to \R^n$. The tangent map of $Y$ is
$T(Y)(x, w) = \bigl(x, \bar{Y}(x), w, d\bar{Y}_x(w)\bigr)$, so
\begin{align*}
  (T(Y) \circ X)(x)
    = \bigl(x,\; \bar{Y}(x),\; \bar{X}(x),\; d\bar{Y}_x(\bar{X}(x))\bigr),
\end{align*}
and applying $K_\Gamma$ yields
\begin{align*}
  (\nabla^{K_\Gamma}_X Y)(x)
    = \bigl(x,\; d\bar{Y}_x(\bar{X}(x))
      + \Gamma(x)(\bar{Y}(x), \bar{X}(x))\bigr).
\end{align*}
For the canonical connection $K_0$ this is the directional derivative
$\nabla^{K_0}_X Y = d\bar{Y}(\bar{X})$, and in general it is the classical
local expression of a covariant derivative through Christoffel symbols.
Note the order of the arguments of $\Gamma(x)$: the second component of
$T(Y) \circ X$ (over $\pi_{TM}$) carries the value of $Y$ and the third
(over $T(\pi_M)$) that of $X$; when $\Gamma(x)$ is symmetric, i.e.\ when
$K_\Gamma$ is torsion-free, the order is immaterial.
\end{ex}

% ----------------------------------------
\subsection{Horizontal Connections, Full Connections, and Effective Connections}%
\label{subsec:full-connection}
% ----------------------------------------

\begin{defn}[{\cite[Definitions~4.5, 4.6]{CC17}, modified}]\label{def:horizontal-connection}
An \textbf{affine horizontal connection} on $M$ is a
smooth morphism $H \colon TM \times_M TM \to T^2M$ satisfying the following conditions
with respect to the projections $\pr_1, \pr_2 \colon TM \times_M TM \to TM$.
We denote the scalar multiplications on the fibre product $TM \times_M TM$
induced by $\kappa_M$ on each factor by
\begin{align*}
  \kappa^{(1)}, \kappa^{(2)} \colon \R \times (TM \times_M TM) \to TM \times_M TM,
\end{align*}
that is, $\kappa^{(1)}(r, (v_0, v_1)) = (\kappa_M(r, v_0), v_1)$ and
$\kappa^{(2)}(r, (v_0, v_1)) = (v_0, \kappa_M(r, v_1))$.
\begin{enumerate}[(a)]
    \item $T(\pi_M) \circ H = \pr_1$
    \item $\pi_{TM} \circ H = \pr_2$
    \item $\lambda_{TM} \circ H = T(H) \circ (\lambda_M \times 0_{TM})$
    \item $\tau_{TM} \circ T(\lambda_M) \circ H = T(H) \circ (0_{TM} \times \lambda_M)$
    \item $H \circ \kappa^{(1)} = \kappa_{TM} \circ (\id_{\R} \times H)$
    \item $H \circ \kappa^{(2)} = T_{(2)}\kappa_M \circ (\id_{\R} \times H)$
\end{enumerate}
\end{defn}

\begin{rem}\label{rem:horizontal-scalar-axiom}
  Axioms~(e) and~(f) express the preservation of scalar multiplication with respect
  to each of the two tangent bundle structures on $T^2M$:
  $\pi_{TM} \colon T^2M \to TM$ (with scalar multiplication $\kappa_{TM}$) and
  $T(\pi_M) \colon T^2M \to TM$ (with scalar multiplication $T_{(2)}\kappa_M$).
\end{rem}

\begin{ex}[{\cite[Example~3.6]{LW}}]%
\label{ex:euclidean-horizontal-connection}
In the coordinates of Example~\ref{ex:euclidean-vertical-connection}, the
\emph{canonical horizontal connection}
$H \colon T\R^n \times_{\R^n} T\R^n \to T^2\R^n$ is given by
\begin{align*}
    H\bigl((u, v_0),\, (u, v_1)\bigr) = (u, v_1, v_0, 0).
\end{align*}

One verifies both scalar axioms directly: $\kappa_{TM}$ scales the $(v_1, v_{01})$ components, which carry the first argument by axiom~(a), while $T_{(2)}\kappa_M$ scales the $(v_0, v_{01})$ components, which carry the second argument by axiom~(b); the last component of $H$ vanishes.
\end{ex}

\begin{defn}[{\cite[Definition~5.2]{CC17}}]\label{def:full-connection}
A pair $(K, H)$ of a vertical connection $K \colon T^2M \to TM$ and a horizontal
connection $H \colon TM \times_M TM \to T^2M$ on $M$
is called a \textbf{full connection} if the
following two conditions hold:
\begin{enumerate}[(a)]
    \item $K \circ H = 0_M \circ \pi_M \circ \pr_1$,
    \item $\mu \circ \langle K, \pi_{TM} \rangle + H \circ U = \id_{T^2M}$.
\end{enumerate}
Here
\begin{align*}
  \mu &\coloneqq T({+}) \circ ((\lambda \circ \pr_1) \times (0 \circ \pr_2))
    \colon TM \times_M TM \to T^2M, \\
  U &\coloneqq \langle T(\pi_M), \pi_{TM} \rangle
    \colon T^2M \to TM \times_M TM;
\end{align*}
condition (b) states that every element of $T^2M$ is the sum of its
vertical and horizontal parts.
\end{defn}

\begin{prop}[{\cite[Proposition~5.8]{CC17}}]\label{prop:full-connection-fibre-product}
  Let $(K, H)$ be a full connection on $M$.
  Then the following is a fibre product diagram:
  \begin{equation}\label{eq:cx_ld}
\xymatrix{
            & T^2M \ar[dl]_{\pi_{TM}} \ar[d]|{T(\pi_M)} \ar[dr]^K & \\
TM \ar[dr]_{\pi_M} & TM \ar[d]|{\pi_M}                          & TM \ar[dl]^{\pi_M}\\
   & M                                                 &
}
\end{equation}
  That is,
  \begin{align*}
    T^2M \cong TM \times_M TM \times_M TM.
  \end{align*}
\end{prop}

\begin{rem}\label{rem:theta-bi-scalar}
Recall that $T^2M$ carries two scalar multiplications over $TM$: the
multiplication $\kappa_{TM} \colon \R \times T^2M \to T^2M$ for the bundle
$\pi_{TM}$, and the multiplication $T_{(2)}\kappa_M \colon \R \times T^2M \to T^2M$
for the bundle $T(\pi_M)$, where
\[
  T_{(2)}\kappa_M = \tau_M \circ \kappa_{TM} \circ (\id_{\R} \times \tau_M)
\]
is the partial tangent of $\kappa_M$ (Definition~\ref{def:scalar-multiplication}).
On the fibre product $TM \times_M TM \times_M TM$ we write $\kappa^{(2,3)}$
(resp.\ $\kappa^{(1,3)}$) for the scalar multiplication that scales the second and
third (resp.\ the first and third) factors by $\kappa_M$ and leaves the remaining
factor fixed.

If the vertical connection $K$ satisfies axioms (e) and (f) of
Definition~\ref{def:vertical-connection}, then the isomorphism
$\Theta = \langle \pi_{TM},\, T(\pi_M),\, K \rangle$ of
Proposition~\ref{prop:full-connection-fibre-product} is moreover an isomorphism of
$\R$-modules for \emph{both} structures:
\begin{align*}
  \Theta \circ \kappa_{TM}
    &= \kappa^{(2,3)} \circ (\id_{\R} \times \Theta),
  &
  \Theta \circ T_{(2)}\kappa_M
    &= \kappa^{(1,3)} \circ (\id_{\R} \times \Theta).
\end{align*}
Indeed, the three components are checked separately.
For $\kappa_{TM}$: the projection $\pi_{TM}$ is left fixed (it is the base of
$\kappa_{TM}$), while
\[
  T(\pi_M) \circ \kappa_{TM} = \kappa_M \circ (\id_{\R} \times T(\pi_M))
  \quad\text{(naturality of $\kappa$),}
  \qquad
  K \circ \kappa_{TM} = \kappa_M \circ (\id_{\R} \times K)
  \quad\text{(axiom (e)).}
\]
For $T_{(2)}\kappa_M$: the projection $T(\pi_M)$ is left fixed, while, using
$\pi_{TM} \circ \tau_M = T(\pi_M)$ together with naturality of $\kappa$,
\[
  \pi_{TM} \circ T_{(2)}\kappa_M = \kappa_M \circ (\id_{\R} \times \pi_{TM}),
  \qquad
  K \circ T_{(2)}\kappa_M = \kappa_M \circ (\id_{\R} \times K)
  \quad\text{(axiom (f)).}
\]
In particular both $\kappa_{TM}$ and $T_{(2)}\kappa_M$ scale the third
($K$-)factor, reflecting that the second-order component of $T^2M$ is scaled by
each of the two structures.
\end{rem}

\begin{defn}[{\cite[Definition~6.6, Theorem~8.4]{LW}}]\label{def:effective-vertical-connection}
An affine vertical connection $K \colon T^2M \to TM$ on $M$ is called
\textbf{effective} if the morphism
\begin{align*}
    \Theta \colon T^2M \longrightarrow TM \times_M TM \times_M TM,
    \qquad W \longmapsto \bigl(\pi_{TM}(W),\; T(\pi_M)(W),\; K(W)\bigr)
\end{align*}
is a diffeomorphism of elastic diffeological spaces.
\end{defn}

\begin{prop}[{\cite[Theorem~7.3]{LW}}]\label{prop:effective-iff-full}
Given an effective vertical connection $K \colon T^2M \to TM$ on $M$,
there exists a unique horizontal connection
$H \colon TM \times_M TM \to T^2M$ compatible with $K$
(i.e., such that $(K, H)$ is a full connection).
Consequently, specifying a full connection is equivalent to choosing an effective vertical connection.
\end{prop}

\begin{rem}\label{rem:H-inherits-scalar}
The construction in \cite[Theorem~7.3]{LW} is carried out for differential bundles,
whose structure involves no scalar multiplication, so the morphism $H$ produced by
Proposition~\ref{prop:effective-iff-full} need not a priori respect the scalar
multiplications. However, the horizontal connection $H$ is precisely
\[
  H = \Theta^{-1} \circ s, \qquad
  s(v_1, v_2) = \bigl(v_1,\; v_2,\; 0_M \circ \pi_M(v_1)\bigr),
\]
that is, the inverse image under $\Theta$ of the inclusion of $TM \times_M TM$ as
the first two factors. Hence, if $K$ additionally satisfies axioms (e) and (f),
then by Remark~\ref{rem:theta-bi-scalar} scaling the $\pi_{TM}$-input of $H$
corresponds to $T_{(2)}\kappa_M$ and scaling the $T(\pi_M)$-input corresponds to
$\kappa_{TM}$; consequently $H$ satisfies the scalar axioms (e) and (f) of
Definition~\ref{def:horizontal-connection}, with no further hypotheses.%
\footnote{We adopt the convention $\pi_{TM} \circ H = \pr_2$ and
$T(\pi_M) \circ H = \pr_1$ in Definition~\ref{def:horizontal-connection}
(consistent with the canonical horizontal connection on $\R^n$,
Example~\ref{ex:euclidean-horizontal-connection}); with this convention
$\kappa^{(2)}$ is paired with $T_{(2)}\kappa_M$ in axiom (f) and $\kappa^{(1)}$
with $\kappa_{TM}$ in axiom (e).}
Thus the equivalence of Proposition~\ref{prop:effective-iff-full} holds at the
level of structures preserving scalar multiplication: an effective vertical
connection satisfying (e), (f) corresponds to a full connection whose horizontal
part satisfies (e), (f).
\end{rem}

% ----------------------------------------
\subsection{Covariant Derivatives along Maps}\label{subsec:cov-along}
% ----------------------------------------

\begin{defn}\label{def:vf-along-map}
Let $f \colon M \to N$ be a smooth map of elastic diffeological spaces.
A smooth map $V \colon M \to TN$ is called a
\textbf{vector field along $f$} if $\pi_N \circ V = f$.
\end{defn}

\begin{defn}\label{def:covariant-derivative-along}
Let $K$ be a vertical connection on $N$, and let $V$ be a vector field
along $f \colon M \to N$.
The \textbf{covariant derivative} of $V$ in the direction of a vector field
$X \colon M \to TM$ on $M$ is defined by
\begin{align*}
    \nabla^K_X V := K \circ T(V) \circ X.
\end{align*}
When $\nabla^K_X V = 0_N \circ f$, we say that $V$ is \textbf{parallel}
in the direction of $X$.
\end{defn}

% ============================================================
\section{Riemannian Metrics and the Levi-Civita Connection}\label{sec:riemannian-metric}
% ============================================================

Throughout this section, $M$ denotes an elastic diffeological space unless
stated otherwise. Tangent vectors $v, w \in TM$ are said to lie \emph{in a
common fibre} if $\pi_M(v) = \pi_M(w)$; each fibre of $\pi_M$ is an
$\R$-module with respect to the operations $+_M$ and $\kappa_M$ of
Theorem~\ref{thm:tangent-structure-2}.

% ----------------------------------------
\subsection{Riemannian Metrics}\label{subsec:metric}
% ----------------------------------------

\begin{defn}\label{def:metric}
A smooth map $g \colon TM \times_M TM \to \R$ is called a
\textbf{Riemannian metric} on $M$ if it satisfies the following conditions
for all $v, w, z \in TM$ lying in a common fibre and all
$r_1, r_2 \in \R$:
 \begin{description}
     \item[Bilinearity]
   $g(r_1v + r_2w,\, z) = r_1g(v, z) + r_2g(w, z)$,\quad
   $g(v,\, r_1w + r_2z) = r_1g(v, w) + r_2g(v, z)$.
 \item[Symmetry]
   $g(v, w) = g(w, v)$ for all $(v, w) \in TM \times_M TM$.
 \item[Positivity]
   $g(v, v) \geq 0$ for all $v \in TM$.
 \item[Definiteness]
   $g(v, v) = 0$ if and only if $v = 0_M(\pi_M(v))$.
 \end{description}
\end{defn}

In \cite{KSS}, weak Riemannian metrics were defined for arbitrary
diffeological spaces as follows.

\begin{defn}[{\cite[Definition~3.1]{KSS}}]\label{def:kss-weak-metric}
  Let $M$ be a diffeological space, not necessarily elastic. A map
  $g \colon \mathbb{L}\hat{T}_2(M) \to \R$ is a
  \textbf{weak Riemannian metric} in the sense of \cite{KSS} if,
  for every $P \in \D(M)$, the map
  \begin{align*}
    g(P) := g \circ j_P \colon TU_P \times_{U_P} TU_P \to \R
  \end{align*}
  is a positive symmetric covariant $2$-tensor field on $U_P$.
  Here $j_P \colon \hat{T}_2U_P = TU_P \times_{U_P} TU_P \to \mathbb{L}\hat{T}_2(M)$
  denotes the canonical map into the colimit
  $\mathbb{L}\hat{T}_2(M) = \colim_{P \in \D(M)} \hat{T}_2U_P$,
  written $\pi_{\hat{T}_2(U_P)}$ in \cite{KSS}.
\end{defn}

\begin{defn}[{\cite[Definition~3.7]{KSS}}]\label{def:kss-definite}
  Let $M$ be a diffeological space and let $g$ be a weak Riemannian metric
  on $M$. If there exists a generating family $\mathcal{G}$ of $\D(M)$ such
  that $g(P)$ is definite in the classical sense for every
  $P \in \mathcal{G}$, then $g$ is said to be \textbf{definite} in the
  sense of \cite{KSS}.
\end{defn}

\begin{rem}\label{rem:kss-equivalence}
  Since $M$ is elastic, the isomorphism
  $\theta_{2,M} \colon \mathbb{L}\hat{T}_2(M) \xrightarrow{\ \simeq\ } TM \times_M TM$
  of axiom~(E1) identifies the two sides; under this identification the
  canonical map $j_P$ corresponds to
  $T(P) \times_{U_P} T(P)$ (Remark~\ref{rem:cocone-TP}).
  Consequently, on an elastic diffeological space the axioms of a weak
  metric in the sense of \cite{KSS} coincide with bilinearity, symmetry,
  and positivity in the sense of Definition~\ref{def:metric}.
\end{rem}

\begin{notation}\label{not:g-along-plot}
  For a plot $P \in \D(M)$ we write
  \[
    g(P) := g \circ \bigl(T(P) \times_P T(P)\bigr),
  \]
  a positive symmetric covariant $2$-tensor field on $U_P$, and
  $g(P)(u)(v, w) := g([P, u, v], [P, u, w])$ for its value at $u \in U_P$.
  This notation recurs in the energy and length functionals of
  Section~\ref{sec:geodesics}.
\end{notation}

\begin{prop}\label{prop:kss-definite-implies-definite}
  Let $g$ be a weak Riemannian metric on $M$ which is definite in the
  sense of \cite{KSS}. Then $g$, regarded as a map
  $TM \times_M TM \to \R$ via Remark~\ref{rem:kss-equivalence},
  satisfies the definiteness condition of Definition~\ref{def:metric}.
\end{prop}

\begin{proof}
  Let $[P, u, v] \in TM$ with $g([P, u, v], [P, u, v]) = 0$.
  Since the generating family $\mathcal{G}$ generates $\D(M)$, there exists
  an open neighbourhood $W \subseteq U_P$ of $u$ such that either $P|_W$ is
  constant, or $P|_W = Q \circ f$ for some generator
  $Q \in \mathcal{G}$ and some smooth map $f \colon W \to U_Q$.
  If $P|_W$ is constant, then $[P, u, v] = 0_M(P(u))$ and there is nothing
  to prove. In the latter case, functoriality
  $T(P|_W) = T(Q) \circ \hat{T}(f)$ (Remark~\ref{rem:cocone-TP}) gives
  \begin{align*}
    [P, u, v] = [P|_W, u, v] = [Q, f(u), df_u(v)].
  \end{align*}
  Hence, in terms of the pullback tensors of
  Definition~\ref{def:kss-weak-metric},
  \begin{align*}
    0 = g([P, u, v], [P, u, v])
      = g(Q)(f(u))\bigl(df_u(v),\, df_u(v)\bigr).
  \end{align*}
  Since $Q \in \mathcal{G}$, the tensor $g(Q)$ is definite, so
  $df_u(v) = 0$. Therefore
  $[P, u, v] = [Q, f(u), 0] = 0_M(P(u))$,
  which is the definiteness condition of Definition~\ref{def:metric}.
\end{proof}

Recall from Lemma~\ref{lem:cotangent-omega} the identification
$\Clin(TM, \R) \cong \Omega^1(M)$. A Riemannian metric induces a
map from vector fields to $1$-forms as follows: for $V \in \mathfrak{X}(M)$,
the composite
\begin{align*}
  g \circ \langle V \circ \pi_M,\, \id_{TM} \rangle \colon TM \to \R,
  \qquad u \longmapsto g\bigl(V(\pi_M(u)),\, u\bigr),
\end{align*}
is smooth and fibrewise linear by bilinearity of $g$, hence an element of
$\Clin(TM,\R) \cong \Omega^1(M)$.

As in classical Riemannian geometry, a metric induces a correspondence between vector fields and differential forms.

\begin{defn}\label{def:flat}
The \textbf{flat morphism} associated with $g$ is the map
\begin{align*}
  \flat \colon \mathfrak{X}(M) \to \Omega^1(M),
  \qquad
  \flat(V) := g \circ \langle V \circ \pi_M,\, \id_{TM} \rangle.
\end{align*}
\end{defn}

\begin{defn}\label{def:non-degenerate}
Let $g \colon TM \times_M TM \to \R$ be a positive, symmetric, and bilinear map.
\begin{enumerate}
    \item $g$ is \textbf{weakly non-degenerate} if for every
      $v \in TM$, the condition $g(v, u) = 0$ for all $u \in TM$ with
      $\pi_M(u) = \pi_M(v)$ implies $v = 0_M(\pi_M(v))$.
    \item $g$ is \textbf{strongly non-degenerate} if the flat morphism
      $\flat \colon \mathfrak{X}(M) \to \Omega^1(M)$ is a diffeomorphism
      with respect to the functional diffeologies. In this case, its
      inverse is written $\sharp \colon \Omega^1(M) \to \mathfrak{X}(M)$.
      (The notation $\flat$, $\sharp$ follows the classical convention.)
\end{enumerate}
\end{defn}

\begin{rem}\label{rem:weak-non-degenerate-linear}
  Weak non-degeneracy states that $g$ separates the points of each fibre:
  by bilinearity, it is equivalent to the condition that for
  $v, w \in TM$ lying in a common fibre, if $g(v, u) = g(w, u)$ for all
  $u \in TM$ with $\pi_M(u) = \pi_M(v)$, then $v = w$ (apply the
  definition to $v - w$). In particular, if $g$ is weakly non-degenerate,
  then $\flat$ is injective.
\end{rem}

\begin{prop}\label{prop:definite-implies-weak-nondegenerate}
  If $g \colon TM \times_M TM \to \R$ satisfies the definiteness condition,
  then $g$ is weakly non-degenerate.
\end{prop}

\begin{proof}
  Let $v \in TM$ satisfy $g(v, u) = 0$ for all $u \in TM$ with
  $\pi_M(u) = \pi_M(v)$.
  Taking $u = v$ gives $g(v, v) = 0$, and definiteness then implies
  $v = 0_M(\pi_M(v))$.
\end{proof}

% ----------------------------------------
\subsection{The Vertical Levi-Civita Connection}\label{subsec:levi-civita}
% ----------------------------------------

Throughout this subsection, $(M, g)$ denotes an \textbf{elastic Riemannian
diffeological space}, that is, an elastic diffeological space $M$ equipped
with a Riemannian metric $g$ in the sense of Definition~\ref{def:metric}.

Before defining the Levi-Civita connection, we fix notation for the
morphisms appearing in the metric compatibility condition. Since the
pullbacks $T_k$ are preserved pointwise by the tangent functor
(Section~\ref{sec:RosickysAxioms}), we identify
\begin{align*}
  T(TM \times_M TM) \;\cong\; T^2M \times_{T(\pi_M)} T^2M,
\end{align*}
and on this space we define
\begin{align*}
  K \times_{\pi} \pi_{TM} \colon (A, B) \longmapsto (K(A),\, \pi_{TM}(B)),
  \qquad
  \pi_{TM} \times_{\pi} K \colon (A, B) \longmapsto (\pi_{TM}(A),\, K(B)),
\end{align*}
both with values in $TM \times_M TM$; the base points match by axiom~(b)
of Definition~\ref{def:vertical-connection} and naturality of $\pi$.
Finally, for the smooth function $g$ on the diffeological space
$TM \times_M TM$ we write
\begin{align*}
  d(g) := \pr_2 \circ T(g) \colon T(TM \times_M TM) \to \R,
\end{align*}
which restricts along every plot $R$ of $TM \times_M TM$ to the classical
exterior derivative $d(g \circ R)$, by the computation in the proof of
Proposition~\ref{prop:pr2-Tf}.

\begin{defn}\label{def:levi-civita}
A vertical connection $K \colon T^2M \to TM$ on $M$ is called a
\textbf{Levi-Civita connection} for $g$ if it satisfies the following two
conditions:
\begin{enumerate}
    \item $K$ is torsion-free, i.e., $K \circ \tau_M = K$.
    \item \textbf{Metric compatibility}:
    \begin{equation}\label{eq:metric-compatibility}
        d(g) = +_{\R} \circ \langle g \circ (K \times_{\pi} \pi_{TM}),\;
               g \circ (\pi_{TM} \times_{\pi} K) \rangle.
    \end{equation}
\end{enumerate}
\end{defn}

\begin{prop}\label{prop:metric-compatibility-vf}
  Let $K$ be a Levi-Civita connection for $g$.
  For any vector fields $X, Y, Z \in \mathfrak{X}(M)$,
  \begin{align*}
    Xg(Y, Z) = g(\nabla^K_X Y, Z) + g(Y, \nabla^K_X Z).
  \end{align*}
\end{prop}

\begin{proof}
  We compose both sides of \eqref{eq:metric-compatibility} with
  $(T(Y)\circ X,\; T(Z) \circ X) = T(\langle Y, Z\rangle) \circ X$,
  where the equality holds under the identification
  $T(TM \times_M TM) \cong T^2M \times_{T(\pi_M)} T^2M$ above.
  For the left-hand side, functoriality of $T$,
  Proposition~\ref{prop:pr2-Tf} applied to the smooth function
  $g(Y, Z) := g \circ \langle Y, Z \rangle \colon M \to \R$, and
  Definition~\ref{def:vf-action} give
  \begin{align*}
    d(g)\circ T(\langle Y, Z\rangle) \circ X
      = \pr_2 \circ T\bigl(g(Y, Z)\bigr) \circ X
      = d\bigl(g(Y, Z)\bigr) \circ X
      = Xg(Y, Z).
  \end{align*}
  For the right-hand side,
  \begin{align*}
    +_{\R} \circ \langle g \circ (K \times_{\pi} \pi_{TM}),\;
    g \circ (\pi_{TM} \times_{\pi} K) \rangle \circ (T(Y)\circ X, T(Z) \circ X)
      &= g(K \circ T(Y)\circ X,\; \pi_{TM} \circ T(Z) \circ X) \\
      &\quad + g(\pi_{TM} \circ T(Y) \circ X,\; K \circ T(Z)\circ X).
  \end{align*}
  Since $\pi$ is a natural transformation,
  $\pi_{TM} \circ T(Z) \circ X = Z \circ \pi_M \circ X = Z$,
  and similarly $\pi_{TM} \circ T(Y) \circ X = Y$.
  Therefore the right-hand side equals
  $g(\nabla^K_X Y, Z) + g(Y, \nabla^K_X Z)$, which proves the claim.
\end{proof}

\begin{prop}\label{prop:vert-hori-compatibility}
Let $(K, H)$ be a full connection on $M$ whose vertical part $K$ satisfies
the metric compatibility condition \eqref{eq:metric-compatibility}. Then
\begin{align*}
    d(g) \circ (H \times H) = 0,
\end{align*}
where $H \times H$ is defined on the fibre product
$(TM \times_M TM) \times_{\pr_1} (TM \times_M TM)$.
\end{prop}

\begin{proof}
We compose both sides of \eqref{eq:metric-compatibility} on the right with $H \times H$; note that axiom~(a) of horizontal connections, $T(\pi_M) \circ H = \pr_1$, guarantees that $H \times H$ takes values in $T^2M \times_{T(\pi_M)} T^2M \cong T(TM \times_M TM)$.
By axiom~(a) of full connections,
\[
K \circ H = 0_M \circ \pi_M \circ \pr_1.
\]
Using further the property $\pi_{TM} \circ H = \pr_2$ of horizontal connections, each component of the right-hand side satisfies
\[
g \circ (K \times_\pi \pi_{TM}) \circ (H \times H) = 0,
\qquad
g \circ (\pi_{TM} \times_\pi K) \circ (H \times H) = 0,
\]
since $g(0_M(x), u) = 0$ for any $u$ in the fibre over $x$, by
bilinearity. Hence
\[
d(g) \circ (H \times H)
= +_{\R} \circ \langle 0,\; 0 \rangle
= 0. \qedhere
\]
\end{proof}

Using the vector field formulation of metric compatibility, a computation
analogous to that in ordinary Riemannian geometry yields the following
Koszul formula.

\begin{prop}[Koszul formula]\label{prop:koszul-formula}
  Let $K$ be a Levi-Civita connection for $g$.
  For any vector fields $X, Y, Z \in \mathfrak{X}(M)$,
  \begin{equation}\label{eq:koszul}
  \begin{aligned}
    2g(\nabla^K_X Y,\; Z)
      &= X\bigl(g(Y,Z)\bigr) + Y\bigl(g(X,Z)\bigr) - Z\bigl(g(X,Y)\bigr) \\
      &\quad + g([X,Y],Z) - g([X,Z],Y) - g([Y,Z],X).
  \end{aligned}
  \end{equation}
\end{prop}

\begin{proof}
  Applying Proposition~\ref{prop:metric-compatibility-vf} three times,
  \begin{align*}
    X\bigl(g(Y,Z)\bigr) &= g(\nabla^K_X Y, Z) + g(Y, \nabla^K_X Z), \\
    Y\bigl(g(X,Z)\bigr) &= g(\nabla^K_Y X, Z) + g(X, \nabla^K_Y Z), \\
    Z\bigl(g(X,Y)\bigr) &= g(\nabla^K_Z X, Y) + g(X, \nabla^K_Z Y).
  \end{align*}
  Adding the first two identities, subtracting the third, and collecting
  terms by bilinearity of $g$, we obtain
  \begin{align*}
    &X\bigl(g(Y,Z)\bigr) + Y\bigl(g(X,Z)\bigr) - Z\bigl(g(X,Y)\bigr) \\
      &\qquad = g(\nabla^K_X Y + \nabla^K_Y X,\; Z)
        + g(\nabla^K_X Z - \nabla^K_Z X,\; Y)
        + g(\nabla^K_Y Z - \nabla^K_Z Y,\; X).
  \end{align*}
  Since $K$ is torsion-free, Proposition~\ref{prop:torsion-free-bracket}
  gives $\nabla^K_X Z - \nabla^K_Z X = [X, Z]$,
  $\nabla^K_Y Z - \nabla^K_Z Y = [Y, Z]$, and
  $\nabla^K_X Y + \nabla^K_Y X = 2\nabla^K_X Y - [X, Y]$. Substituting,
  \begin{align*}
    X\bigl(g(Y,Z)\bigr) + Y\bigl(g(X,Z)\bigr) - Z\bigl(g(X,Y)\bigr)
      = 2g(\nabla^K_X Y, Z) - g([X,Y], Z) + g([X,Z], Y) + g([Y,Z], X),
  \end{align*}
  which rearranges to \eqref{eq:koszul}.
\end{proof}

\begin{thm}\label{thm:uniqueness-cov-deriv}
  If $K_1$ and $K_2$ are Levi-Civita connections for $g$, then
  $\nabla^{K_1} = \nabla^{K_2}$.
\end{thm}

\begin{proof}
  Fix $X, Y \in \mathfrak{X}(M)$ and let
  $D := \nabla^{K_1}_X Y - \nabla^{K_2}_X Y \in \mathfrak{X}(M)$,
  the difference being taken in the fibrewise abelian group structure of
  $TM$. The right-hand side of the Koszul formula \eqref{eq:koszul} does
  not involve $K$, so
  $g(\nabla^{K_1}_X Y, Z) = g(\nabla^{K_2}_X Y, Z)$, and hence by
  bilinearity $g(D, Z) = 0$, for every $Z \in \mathfrak{X}(M)$.
  Since $D$ is itself a vector field, we may substitute $Z = D$.
  This yields $g(D(x), D(x)) = 0$ for every $x \in M$, and
  definiteness gives $D(x) = 0_M(x)$. Therefore
  $\nabla^{K_1}_X Y = \nabla^{K_2}_X Y$.
\end{proof}

\begin{rem}\label{rem:uniqueness-weak-metric}
  The proof above uses only that the vector field $D$ itself may be
  inserted as $Z$; in particular it does not require every tangent vector
  to be realized by a global vector field. Definiteness may moreover be
  relaxed: if $g$ is merely positive, symmetric, and bilinear (a weak
  metric) and weakly non-degenerate, the same conclusion holds. Indeed,
  positivity implies the fibrewise Cauchy--Schwarz inequality
  $g(v, u)^2 \leq g(v, v)\, g(u, u)$ for $u, v \in TM$ lying in a common
  fibre, so $g(D(x), D(x)) = 0$ forces $g(D(x), u) = 0$ for all
  $u \in TM$ with $\pi_M(u) = x$, and weak non-degeneracy gives
  $D(x) = 0_M(x)$.
\end{rem}

\begin{prob}\label{prob:open-questions}
  The following problems remain open.
  \begin{enumerate}
    \item Does the converse of Proposition~\ref{prop:vert-hori-compatibility} hold?
    \item Can Theorem~\ref{thm:uniqueness-cov-deriv} be strengthened to
      uniqueness of the vertical connection itself?
      Also, in analogy with the classical theory, does $K$ exist whenever
      $g$ is strongly non-degenerate?
  \end{enumerate}
\end{prob}

% ============================================================
\section{Curvature}\label{sec:curvature}
% ============================================================

In this section, following \cite[Section~3.16]{CC17}, we introduce the curvature
morphism and curvature tensor on elastic diffeological spaces.

\begin{defn}[{\cite[Definition~3.20]{CC17}}]\label{def:curvature-morphism}
  Let $M$ be an elastic diffeological space and $K \colon T^2M \to TM$ a vertical
  connection on $M$. The \textbf{curvature morphism} $C_K \colon T^3M \to TM$ of $K$
  is defined by
  \begin{align*}
    C_K := K \circ T(K) \circ \tau_{TM} - K \circ T(K).
  \end{align*}
 Both composites $K \circ T(K) \circ \tau_{TM}$ and $K \circ T(K)$ are morphisms $T^3M \to TM$ lying over the same morphism $\pi_M \circ \pi_{TM} \circ \pi_{T^2M} \colon T^3M \to M$, by axiom~(b) of Definition~\ref{def:vertical-connection} and naturality of $\pi$; their difference is therefore formed in the fibrewise abelian group structure of the bundle $\pi_M \colon TM \to M$, whose fibres admit negatives, being $\R$-modules. 
 The equation $C_K = 0$ means that $C_K$ factors through the zero section, i.e.\ $C_K = 0_M \circ \pi_M \circ \pi_{TM} \circ \pi_{T^2M}$; in this case the connection $K$ is said to be \textbf{flat}.
\end{defn}


\begin{ex}\label{ex:flat-euclidean}
  Let $M = \R^n$ with the canonical affine vertical connection
  $K_0(u, v_0, v_1, v_{01}) = (u, v_{01})$ of
  Example~\ref{ex:euclidean-vertical-connection}. Extending the coordinates
  of that example, identify $T^3\R^n \cong (\R^n)^{8}$ with coordinates
  $(v_S)_{S \subseteq \{0,1,2\}}$, where $v_\emptyset = u$ and the
  subscripts record the tangent directions; the first four components
  $(u, v_0, v_1, v_{01})$ form the base point in $T^2\R^n$ and the last
  four $(v_2, v_{02}, v_{12}, v_{012})$ its derivative in the new
  direction $2$. The morphism $\tau_{TM}$, being the component of the
  canonical flip at the object $TM$, exchanges the two outer directions
  $1$ and $2$ while fixing the direction $0$ internal to $TM$:
  \begin{align*}
    \tau_{TM}(v_S)_S = (v_{\sigma(S)})_S, \qquad \sigma = (1\ 2).
  \end{align*}
  Since $K_0$ is linear, its tangent map acts by $K_0$ on both blocks:
  \begin{align*}
    T(K_0)(v_S)_S = (u,\, v_{01},\, v_2,\, v_{012}) \in T^2\R^n,
  \end{align*}
  whence $K_0 \circ T(K_0) = (u, v_{012})$. Precomposing with
  $\tau_{TM}$ replaces $(v_1, v_{01})$ by $(v_2, v_{02})$ and vice versa,
  so
  \begin{align*}
    K_0 \circ T(K_0) \circ \tau_{TM}
      = K_0(u,\, v_{02},\, v_1,\, v_{012}) = (u, v_{012}).
  \end{align*}
  The two composites coincide, and therefore $C_{K_0} = 0$: the canonical
  connection on $\R^n$ is flat.
\end{ex}

\begin{defn}[{\cite[Definition~3.21]{CC17}}]\label{def:curvature-tensor}
  Let $X, Y, Z \in \mathfrak{X}(M)$ be vector fields on $M$.
  The \textbf{curvature tensor} $R^K(X,Y,Z) \in \mathfrak{X}(M)$
  is defined using the curvature morphism $C_K$ by
  \begin{align*}
    R^K(X,Y,Z) := C_K \circ T^2(Z) \circ T(X) \circ Y.
  \end{align*}
\end{defn}


This curvature tensor admits an expression in terms of covariant derivatives,
as in ordinary Riemannian geometry.

\begin{prop}[{\cite[Proposition~3.22]{CC17}}]\label{prop:curvature-covariant-derivative}
  For vector fields $X, Y, Z \in \mathfrak{X}(M)$ on $M$, the following holds:
  \begin{align*}
    R^K(X,Y,Z)
      = \nabla^K_X (\nabla^K_Y Z)
        - \nabla^K_Y (\nabla^K_X Z)
        - \nabla^K_{[X,Y]} Z.
  \end{align*}
\end{prop}


\begin{rem}\label{rem:curvature-pending}
  The curvature recalled above enjoys the familiar properties of the classical theory, established by Cockett and Cruttwell at the level of an arbitrary tangent category: the expression of $R^K$ in terms of the covariant derivative (Proposition~\ref{prop:curvature-covariant-derivative} above), the identities for curvature and torsion, and, for torsion-free connections, the Bianchi identities \cite[Theorem 3.34]{CC17}. These results require the tangent
  bundles to admit negatives, which holds in $\mathsf{Elast}$ since the fibres of $T(M)$ are $\R$-modules; they therefore apply verbatim in the present setting, and we do not reproduce the proofs. 

  In this paper we do not pursue the interaction between curvature and a Riemannian metric; the metric symmetries of the curvature tensor and sectional curvature on elastic Riemannian diffeological spaces are left to future work.
\end{rem}

% ============================================================
\section{Geodesics}\label{sec:geodesics}
% ============================================================

Throughout this section, $M$ denotes an elastic diffeological space, and
$(K, H)$ a full connection on $M$ (Definition~\ref{def:full-connection}),
unless further hypotheses are stated. We write $\Path(M)$ for the
diffeological space $C^\infty(\R, M)$ of smooth paths, with the functional
diffeology, and for $x, y \in M$ we write $\Path(M; x, y)$ for the
subspace of paths $\gamma$ with $\gamma(0) = x$ and $\gamma(1) = y$.

% ----------------------------------------
\subsection{Dynamical Systems and Curve Objects}\label{subsec:parallel-transport}
% ----------------------------------------

Unlike in ordinary Riemannian geometry, the parallel transport of a given tangent
vector $v \in TN$ along a smooth map $f$ does not always exist globally.
To accommodate parallel transport, we introduce the following notions of
dynamical system and curve object.

\begin{defn}[{\cite[Definition~2.24]{CCL}}]\label{def:dynamical-system}
  Let $V$ be a vector field on $M$ and $x \in M$.
  The triple $(M, V, x)$ is called a \textbf{dynamical system};
  $V$ is the \textbf{differential transition} and $x$ is the \textbf{initial state}.
  A \textbf{morphism} $f \colon (M, V, x) \to (M', V', x')$ of dynamical systems
  is a smooth map $f \colon M \to M'$ satisfying $f(x) = x'$ and
  $T(f) \circ V = V' \circ f$; that is, the following diagram commutes:
  \[ \xymatrix{1 \ar[r]^x \ar[dr]_{x'} & M \ar[d]_f \ar[r]^V & T(M) \ar[d]^{T(f)} \\
               & M' \ar[r]_{V'} & T(M') } \]
\end{defn}

\begin{rem}\label{rem:initial-state}
  In a general tangent category, the initial state is a morphism
  $x \colon \mathbf{1} \to M$ from the terminal object $\mathbf{1}$
  (\cite[Definition~2.24]{CCL}).
  In elastic diffeology, since each object is a set, one may take points directly;
  however, individual points can also be viewed as morphisms from $\mathbf{1}$
  when convenient.
\end{rem}

\begin{defn}[{\cite[Definition~3.1]{CCL}}]\label{def:parametrized-dynamical-system}
  Let $V$ be a vector field on $M$ and $g \colon X \to M$ a smooth map.
  The triple $(M, V, g)$ is called a \textbf{parametrized dynamical system},
  and $g$ is the \textbf{parametrized initial state}.
\end{defn}

\begin{defn}[{\cite[Definition~3.1]{CCL}}]\label{def:parametrized-solution}
  Let $C$ be an elastic diffeological space, $c_1 \colon C \to TC$ a vector field
  on $C$, and $c_0 \in C$. A \textbf{solution} to the parametrized dynamical system
  $(M, V, g)$ is a smooth map $\gamma \colon C \times X \to M$ satisfying
  the initial condition $\gamma(c_0, p) = g(p)$ and the differential condition
  $T(\gamma) \circ (c_1 \times 0_X) = V \circ \gamma$.
  That is, the following diagram commutes:
 \[
   \xymatrix{
     X \ar[r]^{\langle c_0, 1\rangle} \ar[dr]_{g}
       & C \times X \ar[r]^{c_1 \times 0} \ar[d]_{\gamma}
       & T(C \times X) \ar[d]^{T(\gamma)} \\
     & M \ar[r]_{V}
       & TM
   }
 \]
\end{defn}

\begin{rem}\label{rem:point-as-constant-map}
  As in Remark~\ref{rem:initial-state}, $c_0 \in C$ may also be regarded, for
  any object $X$, as the constant map $X \to C$ sending every point to $c_0$;
  we write this map simply as $c_0$ as well.
\end{rem}

\begin{defn}[{\cite[Definition~4.1]{CCL}}]\label{def:curve-object}
  A dynamical system $(C, c_1, c_0)$ is a \textbf{curve object} if it satisfies
  the following:
  \begin{enumerate}
    \item \textbf{Preinitiality}: For any parametrized dynamical system $(M, V, g)$,
      a solution $\gamma \colon C \times X \to M$, if it exists, is unique:
    \[
      \xymatrix{
        X \ar[r]^{\langle c_0, 1\rangle} \ar[dr]_{g}
          & C \times X \ar[r]^{c_1 \times 0} \ar@{-->}[d]_{\gamma}
          & T(C \times X) \ar@{-->}[d]^{T(\gamma)} \\
        & M \ar[r]_{V}
          & TM
      }
    \]
    \item \textbf{Self-commutativity}: $\tau_C \circ T(c_1) \circ c_1
      = T(c_1) \circ c_1$.
    \item \textbf{Completeness}: The parametrized dynamical system
      $(C, c_1, \id_C)$ on $C$ has a unique solution $\sigma \colon C \times C \to C$:
    \[
      \xymatrix{
        C \ar[r]^{\langle c_0, 1\rangle} \ar[dr]_{\id_C}
          & C \times C \ar[r]^{c_1 \times 0} \ar@{-->}[d]_{\sigma}
          & T(C \times C) \ar@{-->}[d]^{T(\sigma)} \\
        & C \ar[r]_{c_1}
          & TC
      }
    \]
  \end{enumerate}
\end{defn}

\begin{defn}[{\cite[Definition~4.9]{CCL}}]\label{def:complete-vector-field}
  A vector field $V$ on $M$ is called \textbf{complete} if the parametrized
  dynamical system $(M, V, \id_M)$ with parametrized initial state $\id_M$
  has a solution $\gamma \colon C \times M \to M$:
  \[
    \xymatrix{
      M \ar[r]^{\langle c_0, 1\rangle} \ar[dr]_{\id_M}
        & C \times M \ar[r]^{c_1 \times 0} \ar@{-->}[d]_{\gamma}
        & T(C \times M) \ar@{-->}[d]^{T(\gamma)} \\
      & M \ar[r]_{V}
        & TM
    }
  \]
\end{defn}

A curve object plays the role of $\R$ (or an interval therein) in a general tangent
category. In the differential geometry of finite-dimensional manifolds,
parallel transport of vector fields along paths reduces to a linear ODE in local
coordinates, and the existence and uniqueness of local solutions is guaranteed by
the fundamental theorem of ODEs. In the category of convenient vector spaces,
however, the spaces lack a complete normed structure, so even for smooth vector
fields neither existence nor uniqueness of solutions is guaranteed
(\cite[Remark~4.5]{CCL}). Consequently, solutions to dynamical systems need not
exist even when the domain is $\R$, and hence parallel transport need not always
exist. The same applies in the category of elastic diffeological spaces,
which contains the category of convenient vector spaces.

Parallel transport always exists along curve objects satisfying the condition
of \emph{linear completeness}. See \cite[Theorem~5.20]{CC17} and
\cite[Remark~5.4]{CCL} for details.

% ----------------------------------------
\subsection{Definition of Geodesics}\label{subsec:geodesic-def}
% ----------------------------------------

Let $\gamma$ be a path on $M$, and let $\partial$ denote the unit vector field
on $\R$. We call $\dot{\gamma} = T(\gamma) \circ \partial = [\gamma, t, \partial]$
the \textbf{velocity vector field} of $\gamma$.

\begin{rem}\label{rem:velocity-along}
  $\dot{\gamma}$ is a vector field along $\gamma$
  (Definition~\ref{def:vf-along-map}). Indeed, by naturality of $\pi$,
  \begin{align*}
    \pi_{M} \circ \dot{\gamma}
      = \pi_{M} \circ T(\gamma) \circ \partial
      = \gamma \circ \pi_{\R} \circ \partial
      = \gamma \circ \id_{\R}
      = \gamma.
  \end{align*}
\end{rem}

\begin{defn}\label{def:geodesic}
Let $K$ be a vertical connection on $M$.
A path $\gamma \colon \R \to M$ is called a \textbf{geodesic} for $K$ if
\begin{align*}
   \nabla^K_{\partial} \dot{\gamma} = 0_M \circ \gamma,
\end{align*}
where $\nabla^K_\partial \dot\gamma = K \circ T(\dot\gamma) \circ \partial$ is
the covariant derivative of the vector field $\dot\gamma$ along $\gamma$
(Definition~\ref{def:covariant-derivative-along}).
\end{defn}

% ----------------------------------------
\subsection{Geodesic Flow}\label{subsec:geodesic-flow}
% ----------------------------------------

\begin{defn}[{\cite[Definition~3.10]{CCL}}]\label{def:nth-order-dynamical-system}
  Let $n \geq 1$ be an integer.
  An \textbf{$n$th-order dynamical system} on $M$ is a triple
  $(M, V, g)$ consisting of an initial state $g \colon X \to T^{n-1}M$ and
  an $n$th-order vector field, i.e., a smooth map
  $V \colon T^{n-1}M \to T^nM$, satisfying
  \begin{align*}
      \pi_{T^{n-1}M} \circ V
        = T(\pi_{T^{n-2}M}) \circ V
        = \cdots
        = T^{n-1}(\pi_M) \circ V
        = \id_{T^{n-1}M}.
  \end{align*}
  When $n = 1$ this coincides with the dynamical system of
  Definition~\ref{def:dynamical-system}.
  Moreover, if $(M, V, g)$ is an $n$th-order dynamical system,
  then $(TM, V, g)$ is an $(n-1)$th-order dynamical system.
\end{defn}

\begin{defn}[{\cite[Definition~3.11]{CCL}}]\label{def:nth-derivative}
  Let $C$ be a curve object.
  For any smooth map $F \colon C \times X \to M$,
  the \textbf{$n$th derivative} $F^{(n)} \colon C \times X \to T^nM$ is
  defined inductively by $F^{(0)} := F$ and, for each $n \geq 1$,
  \begin{align*}
      F^{(n)} := T(F^{(n-1)}) \circ (c_1 \times 0_X).
  \end{align*}
\end{defn}

\begin{rem}\label{rem:first-derivative}
  $F^{(1)}$ is nothing other than $T(F) \circ (c_1 \times 0_X)$,
  the left-hand side of the differential condition
  in the definition of a solution to a first-order dynamical system.
\end{rem}

\begin{defn}[{\cite[Definition~3.12]{CCL}}]\label{def:nth-solution}
  A \textbf{solution} to the $n$th-order dynamical system $(M, V, g)$
  is a smooth map $F \colon C \times X \to M$ satisfying the
  \textbf{initial condition}
  $F^{(n-1)} \circ \langle c_0, \id_X \rangle = g$ and the
  \textbf{differential condition} $F^{(n)} = V \circ F^{(n-1)}$.
\end{defn}

\begin{prop}[{\cite[Proposition~3.13]{CCL}}]\label{prop:higher-dynamical-system}
  Let $(M, V, g)$ be an $n$th-order dynamical system.
  If $\tilde{F} \colon C \times X \to TM$ is a solution to the
  $(n-1)$th-order dynamical system $(TM, V, g)$,
  then $F := \pi_M \circ \tilde{F}$ is a solution to the
  $n$th-order dynamical system $(M, V, g)$. In particular
  $\pi_M \circ F^{(1)} = F$.
\end{prop}

\begin{rem}[{\cite[Corollary~3.14]{CCL}}]\label{rem:reduction-to-first-order}
  By applying the above proposition repeatedly, the problem of solving any
  $n$th-order dynamical system $(M, V, g)$ reduces to that of solving
  the first-order dynamical system $(T^{n-1}M, V, g)$ on the iterated tangent bundle.
\end{rem}

\begin{defn}[{\cite[Section~3.5]{CCL}}]\label{def:geodesic-spray}
  Let $(K, H)$ be a full connection on $M$.
  Define the map $S \colon TM \to T^2M$ by
  \begin{align*}
      S := H \circ \langle \id_{TM}, \id_{TM} \rangle.
  \end{align*}
  Since $S$ satisfies $\pi_{TM} \circ S = \id_{TM}$ and $\pi_{TM} \circ S = T(\pi_M) \circ S = \id_{TM}$,
  it is a second-order vector field on $M$;
  we call it the \textbf{geodesic spray} of $(K, H)$.

  A solution $\Phi \colon C \times TM \to M$ to the second-order dynamical system
  $(M, S, \id_{TM})$ with initial state $\id_{TM} \colon TM \to TM$
  is called the \textbf{geodesic flow} associated with $(K, H)$.
\end{defn}

By definition, the geodesic flow satisfies $\Phi^{(2)} = S \circ \Phi^{(1)}$;
equivalently, the following diagram commutes:
\[
  \xymatrix{
    C \times TM \ar[d]_{\Phi^{(1)}} \ar[dr]^{\Phi^{(2)}} & \\
    TM \ar[r]_{S} & T^2M
  }
\]

\begin{prop}[{\cite[Section~3.5]{CCL}}]\label{prop:acceleration-zero}
  The geodesic flow $\Phi \colon C \times TM \to M$ satisfies
  \begin{align*}
      K \circ \Phi^{(2)} = 0_{M} \circ \Phi.
  \end{align*}
\end{prop}

\begin{proof}
  By the differential condition and Definition~\ref{def:geodesic-spray},
  \begin{align*}
      K \circ \Phi^{(2)}
        = K \circ S \circ \Phi^{(1)}
        = K \circ H \circ \langle \id_{TM}, \id_{TM} \rangle \circ \Phi^{(1)}.
  \end{align*}
  By axiom~(a) of full connections,
  $K \circ H = 0_M \circ \pi_M \circ \pr_1$, so
  \begin{align*}
      K \circ \Phi^{(2)}
        = 0_M \circ \pi_M \circ \pr_1
          \circ \langle \id_{TM}, \id_{TM} \rangle \circ \Phi^{(1)}
        = 0_M \circ \pi_M \circ \Phi^{(1)}
        = 0_M \circ \Phi,
  \end{align*}
  the last equality by $\pi_M \circ \Phi^{(1)} = \Phi$
  (Proposition~\ref{prop:higher-dynamical-system}).
\end{proof}

% ----------------------------------------
\subsection{Existence of Geodesics}\label{subsec:geodesic-existence}
% ----------------------------------------

The subcategory on which $\R$ serves as a curve object plays a distinguished
role in what follows; we name it.

\begin{defn}\label{def:elast-lc}
  Let $\mathsf{Elast}_{\mathrm{curve}}$ denote the full subcategory of $\mathsf{Elast}$ consisting of those elastic diffeological spaces $M$ such that, for every $n \geq 0$, the dynamical system $(\R, \partial, 0)$
  is preinitial for the iterated tangent bundle $T^nM$ (Definition~\ref{def:curve-object}); that is, every parametrized dynamical system on $T^nM$ with domain $\R$ admits at most one solution.
\end{defn}

\begin{rem}\label{rem:elast-curve-closure}
  Conditions (2) and (3) of Definition~\ref{def:curve-object} concern $(\R, \partial, 0)$ itself and hold automatically, the flow of $\partial$ being translation; thus $M \in \mathsf{Elast}_{\mathrm{curve}}$ precisely when $(\R, \partial, 0)$ is a curve object for every $T^nM$. Requiring preinitiality for all iterated tangent bundles, rather than for $M$ alone, makes $\mathsf{Elast}_{\mathrm{curve}}$ stable under the tangent functor---in analogy with the elasticity axiom (E5)---and is exactly what the reduction of higher-order dynamical systems to first-order systems on iterated tangent bundles requires (Remark~\ref{rem:reduction-to-first-order}).
\end{rem}

\begin{rem}\label{rem:manifolds-curve}
  In the category of smooth manifolds, $(\R, \partial, 0)$ is a curve object \cite[Example~4.2]{CCL}. 
  It follows that every smooth manifold belongs to $\mathsf{Elast}_{\mathrm{curve}}$: given two solutions $\gamma, \gamma' \colon \R \times X \to M$ of a parametrized dynamical system on a manifold $M$ with an arbitrary elastic parameter object $X$, restricting along each point $x \colon * \to X$ (using the naturality of the zero section) exhibits $\gamma(-, x)$ and
  $\gamma'(-, x)$ as integral curves of the same vector field with the same initial value; these coincide by preinitiality in the category of manifolds, whence $\gamma = \gamma'$.
\end{rem}

As in ordinary Riemannian geometry,
the existence of geodesics is guaranteed by the geodesic flow.

\begin{thm}\label{thm:geodesic-existence}
  Let $M \in \mathsf{Elast}_{\mathrm{curve}}$ and let $(K, H)$ be a full connection
  on $M$ whose geodesic spray
  $S = H \circ \langle \id_{TM}, \id_{TM} \rangle$ is complete.
  Then for every $v_0 \in TM$ there exists a unique
  geodesic $\gamma \colon \R \to M$ for $K$ satisfying
  \begin{align*}
      \nabla^K_{\partial} \dot{\gamma} = 0_{M} \circ \gamma,
      \qquad
      \dot{\gamma}(0) = v_0.
  \end{align*}
\end{thm}

\begin{proof}
  Since $S$ is complete, the second-order dynamical
  system $(M, S, \id_{TM})$ has a geodesic flow
  $\Phi \colon \R \times TM \to M$ satisfying
  $\Phi^{(2)} = S \circ \Phi^{(1)}$.

  For the given $v_0 \in TM$, let $\iota_{v_0} \colon \R \to \R \times TM$
  be the slice map $t \mapsto (t, v_0)$, and define
  $\gamma := \Phi \circ \iota_{v_0} \colon \R \to M$.
  By functoriality of $T$,
  $\gamma^{(1)} = T(\Phi) \circ T(\iota_{v_0}) \circ \partial$.
  Decomposing $\iota_{v_0} = \langle \id_\R, c_{v_0} \rangle$,
  the product-preservation of $T$ and the fact that
  $T(c_{v_0}) \circ \partial = 0_{TM} \circ c_{v_0}$ give
  \begin{align*}
    T(\iota_{v_0}) \circ \partial
      = \langle \partial,\; 0_{TM} \circ c_{v_0} \rangle
      = (\partial \times 0_{TM}) \circ \iota_{v_0}.
  \end{align*}
  Hence $\gamma^{(1)} = \Phi^{(1)} \circ \iota_{v_0}$,
  and an analogous argument gives $\gamma^{(2)} = \Phi^{(2)} \circ \iota_{v_0}$.
  Composing $\Phi^{(2)} = S \circ \Phi^{(1)}$ on the right with $\iota_{v_0}$
  yields $\gamma^{(2)} = S \circ \gamma^{(1)}$, i.e.,
  $T(\dot{\gamma}) \circ \partial = S \circ \dot{\gamma}$.
  Composing on the left with $K$ and applying
  Proposition~\ref{prop:acceleration-zero} (equivalently
  $K \circ H = 0_M \circ \pi_M \circ \pr_1$),
  the right-hand side becomes
  $0_M \circ \pi_M \circ \dot{\gamma} = 0_M \circ \gamma$,
  which is the geodesic condition $\nabla^K_\partial \dot{\gamma} = 0_M \circ \gamma$.

  Since $M \in \mathsf{Elast}_{\mathrm{curve}}$, i.e.\ $\R$ is a curve object for
  $TM$ and uniqueness of $\Phi$, $\gamma$ with $\dot{\gamma}(0) = v_0$ is uniquely determined.
\end{proof}

% ----------------------------------------
\subsection{The Energy Functional and the First Variation Formula}%
\label{subsec:energy}
% ----------------------------------------

Let $(M, g)$ be an elastic Riemannian diffeological space.
We define the \textbf{energy functional}
$E \colon \Path(M; x, y) \to \R$ by
\begin{align}\label{eq:energy}
  E(\gamma)
    = \frac{1}{2}\int_0^1 g([\gamma, t, \partial], [\gamma, t, \partial]) \, dt
    = \frac{1}{2}\int_0^1 g(\gamma)(t)(\partial, \partial) \, dt,
\end{align}
in the notation of Notation~\ref{not:g-along-plot}; the value $E(\gamma)$ is
finite since $[0,1]$ is compact and the integrand is smooth in $t$.

\begin{defn}[Variation plot]\label{def:variation-plot}
  Let $\gamma \in \Path(M; x, y)$.
  A $1$-plot $P \colon (-\varepsilon, \varepsilon) \to \Path(M; x, y)$
  is called a \textbf{variation plot} of $\gamma$ if it satisfies
  $P(0) = \gamma$, and $P(s)(0) = x$ and $P(s)(1) = y$ for all
  $s \in (-\varepsilon, \varepsilon)$.
\end{defn}

\begin{defn}[Critical path]\label{def:critical-path}
  A path $\gamma \in \Path(M; x, y)$ is called a \textbf{critical path}
  if, for every variation plot $P$ of $\gamma$,
  \begin{align*}
    \left. \frac{d(E \circ P)}{ds} \right|_{s = 0} = 0.
  \end{align*}
\end{defn}

We first record a form of metric compatibility for covariant derivatives along
maps (Definition~\ref{def:covariant-derivative-along}), which is what the
variation computation actually requires: the fields $V_s, V_t$ below are
vector fields along $\ad(P)$, not vector fields on $M$.

\begin{prop}[Metric compatibility along maps]\label{prop:metric-compat-along}
  Let $(M, g)$ be an elastic Riemannian diffeological space with Levi-Civita
  connection $K$, let $f \colon N \to M$ be a smooth map of elastic
  diffeological spaces, and let $V, W$ be vector fields along $f$
  (Definition~\ref{def:vf-along-map}). Then for every $X \in \mathfrak{X}(N)$,
  \begin{align*}
    X\bigl(g(V, W)\bigr)
      = g\bigl(\nabla^K_X V,\, W\bigr) + g\bigl(V,\, \nabla^K_X W\bigr),
  \end{align*}
  where $g(V, W) \colon N \to \R$ is $p \mapsto g(V(p), W(p))$ and
  $\nabla^K_X V := K \circ T(V) \circ X$.
\end{prop}

\begin{proof}
  Since $\pi_M \circ V = f = \pi_M \circ W$, the pair $\langle V, W \rangle$ is
  a smooth map $N \to T_2M = TM \times_M TM$, and $g(V, W) = g \circ \langle V,
  W \rangle$. As $T$ preserves the fibre product $T_2M$, the pair $(T(V) \circ
  X, T(W) \circ X)$ equals $T(\langle V, W \rangle) \circ X$, a map
  $N \to T(T_2M)$. Composing the metric compatibility identity
  \eqref{eq:metric-compatibility} with $T(\langle V, W\rangle) \circ X$ and
  arguing exactly as in Proposition~\ref{prop:metric-compatibility-vf} --- using
  $\pi_{TM} \circ T(V) \circ X = V \circ \pi_N \circ X = V$ and likewise for $W$
  --- gives
  \begin{align*}
    d\bigl(g(V,W)\bigr) \circ X
      = g(K \circ T(V) \circ X,\, W) + g(V,\, K \circ T(W) \circ X),
  \end{align*}
  which is the assertion, since $d(g(V,W)) \circ X = X(g(V,W))$
  (Proposition~\ref{prop:pr2-Tf}, Definition~\ref{def:vf-action}).
\end{proof}

The proof of the first variation formula uses the following symmetry lemma.

\begin{lem}\label{lem:symmetry}
  Let $K$ be a torsion-free vertical connection on $M$.
  For any plot $P \colon U_P \subseteq \R^2 \to M$,
  \begin{align*}
    \nabla^{K}_{\partial_s}[P, (s, t), \partial_t]
      = \nabla^{K}_{\partial_t}[P, (s, t), \partial_s].
  \end{align*}
\end{lem}

\begin{proof}
  By the definition of the covariant derivative along $P$,
  \begin{align*}
    \nabla^{K}_{\partial_s} [P, (s, t), \partial_t]
      = K \circ T^2(P) \circ T(\partial_t) \circ \partial_s,
    \qquad
    \nabla^{K}_{\partial_t}[P, (s, t), \partial_s]
      = K \circ T^2(P) \circ T(\partial_s) \circ \partial_t.
  \end{align*}
  In the coordinates of Example~\ref{ex:euclidean-vertical-connection}, the coordinate fields on $U_P \subseteq \R^2$ have constant components, so that $T(\partial_t)(u, v) = (u, e_2, v, 0)$, and hence
  \begin{align*}
    T(\partial_t) \circ \partial_s(u)
      = (u, e_2, e_1, 0)
      = \tau_{\R^2}(u, e_1, e_2, 0)
      = \tau_{\R^2} \circ T(\partial_s) \circ \partial_t(u),
  \end{align*}
  an instance of the classical relation between the canonical flip and the Lie bracket of vector fields (cf.\ \cite[Lemma~6.13]{KMS}).
  By naturality of $\tau$, $T^2(P) \circ \tau_{\R^2} = \tau_M \circ T^2(P)$, and since $K$ is torsion-free ($K \circ \tau_M = K$),
  \begin{align*}
    \nabla^{K}_{\partial_s}[P, (s, t), \partial_t]
      &= K \circ \tau_M \circ T^2(P) \circ T(\partial_s) \circ \partial_t \\
      &= K \circ T^2(P) \circ T(\partial_s) \circ \partial_t
       = \nabla^{K}_{\partial_t}[P, (s, t), \partial_s]. \qedhere
  \end{align*}
\end{proof}

\begin{lem}[First Variation Formula]\label{lem:first-variation-formula}
  Let $(M, g)$ be an elastic Riemannian diffeological space, and let $K$ be a Levi-Civita connection for $g$. 
  For $\gamma \in \Path(M; x, y)$ and any variation plot $P$ of $\gamma$, consider the vector fields along $\ad(P)$,
  \begin{align*}
    V_s := [\ad(P), (s, t), \partial_s],
    \qquad
    V_t := [\ad(P), (s, t), \partial_t].
  \end{align*}
  Then
  \begin{align*}
      \left. \frac{d(E \circ P)}{ds} \right|_{s = 0}
        = - \int_0^1 g\bigl(V_s,\, \nabla^{K}_{\partial_t} V_t\bigr)
          \Big|_{s=0} \, dt.
  \end{align*}
\end{lem}

\begin{rem}\label{rem:variation-velocity}
For any variation plot $P$, $P(s)$ is a path on $M$ and hence a plot.
For the inclusion $\iota_s \colon \R \to \R^2$, $t \mapsto (s, t)$,
we have $\ad(P) \circ \iota_s = P(s)$. By the colimit description of the
tangent space (Remark~\ref{rem:tangent-structure-2}),
$\bigl[P(s), t, \tfrac{d}{dt}\bigr] = \bigl[\ad(P), (s,t), \partial_t\bigr]$,
so $V_t|_{s}$ is the velocity field of the path $P(s)$.
\end{rem}

\begin{proof}[Proof of Lemma~\ref{lem:first-variation-formula}]
  By \eqref{eq:energy} and differentiation under the integral sign,
  \begin{align*}
      \frac{d(E \circ P)}{ds}
        = \frac{1}{2} \int_0^1 \frac{\partial}{\partial s} g(V_t, V_t) \, dt.
  \end{align*}
  Since $K$ is a Levi-Civita connection and $V_t$ is a vector field along
  $\ad(P)$, metric compatibility along maps
  (Proposition~\ref{prop:metric-compat-along}, with $f = \ad(P)$ and
  $X = \partial_s$) gives
  $\frac{\partial}{\partial s} g(V_t, V_t) = 2g(\nabla^{K}_{\partial_s} V_t, V_t)$,
  hence
  \begin{align*}
      \frac{d(E \circ P)}{ds}
        = \int_0^1 g(\nabla^{K}_{\partial_s} V_t, V_t) \, dt.
  \end{align*}
  By Lemma~\ref{lem:symmetry},
  $\nabla^{K}_{\partial_s} V_t = \nabla^{K}_{\partial_t} V_s$, so applying
  Proposition~\ref{prop:metric-compat-along} once more (now with
  $X = \partial_t$),
  $g(\nabla^{K}_{\partial_t} V_s, V_t)
  = \frac{\partial}{\partial t} g(V_s, V_t)
  - g(V_s, \nabla^{K}_{\partial_t} V_t)$, we obtain
  \begin{align*}
      \frac{d(E \circ P)}{ds}
        = \left[ g(V_s, V_t) \right]_{t=0}^{t=1}
           - \int_0^1 g\bigl(V_s,\, \nabla^{K}_{\partial_t} V_t\bigr) \, dt.
  \end{align*}
  Since the variation plot is constant at the endpoints,
  $V_s|_{t=0} = V_s|_{t=1} = 0$, so the boundary term vanishes.
  Evaluating at $s = 0$ gives the formula.
\end{proof}

\begin{prop}\label{prop:geodesic-critical-path}
  Let $(M, g)$ be as in Lemma~\ref{lem:first-variation-formula}.
  If $\gamma \in \Path(M; x, y)$ is a geodesic for the Levi-Civita connection
  $K$, then $\gamma$ is a critical path of the energy functional.
\end{prop}

\begin{proof}
  For any variation plot $P$ with $P(0) = \gamma$, the field $V_t|_{s=0}$ is the
  velocity $\dot\gamma$ (Remark~\ref{rem:variation-velocity}), so
  $\nabla^{K}_{\partial_t} V_t|_{s=0} = \nabla^K_\partial \dot\gamma
  = 0_M \circ \gamma$ since $\gamma$ is a geodesic. The first variation formula
  (Lemma~\ref{lem:first-variation-formula}) then gives
  $\frac{d(E \circ P)}{ds}\big|_{s=0} = 0$, so $\gamma$ is a critical path.
\end{proof}

Throughout the rest of this subsection we make the following standing
assumptions: $(M, g)$ is an elastic Riemannian diffeological space with
$M \in \mathsf{Elast}_{\mathrm{curve}}$; $(K, H)$ is a full connection on $M$
whose vertical part $K$ is a Levi-Civita connection for $g$; the geodesic
spray $S = H \circ \langle \id_{TM}, \id_{TM} \rangle$ is complete, and
$\Phi \colon \R \times TM \to M$ denotes the geodesic flow of
Theorem~\ref{thm:geodesic-existence}.

\begin{lem}\label{lem:flow-initial}
  For every $v \in TM$,
  \begin{align*}
    \mathrm{(i)}\quad \Phi(0, v) = \pi_M(v),
    \qquad\qquad
    \mathrm{(ii)}\quad \Phi^{(1)}(0, v) = v,
  \end{align*}
  where
  $\Phi^{(1)} = T(\Phi) \circ (\partial \times 0_{TM}) \colon \R \times TM \to TM$
  (Definition~\ref{def:nth-derivative}). In particular, for each $v$ the curve
  $s \mapsto \Phi(s, v)$ is the geodesic with base point $\pi_M(v)$ and initial
  velocity $v$.
\end{lem}

\begin{proof}
  The initial condition of the second-order dynamical system
  $(M, S, \id_{TM})$ (Definition~\ref{def:nth-solution}) is exactly
  $\Phi^{(1)}(0, v) = v$, which is (ii). For (i),
  naturality of $\pi$ gives
  $\pi_M \circ T(\Phi) = \Phi \circ \pi_{\R \times TM}$, and since $T$ preserves
  products (Proposition~\ref{prop:tangent-structure-1}),
  $\pi_{\R \times TM} = \pi_\R \times \pi_{TM}$. Hence
  \begin{align*}
    \pi_M \circ \Phi^{(1)}
      = \Phi \circ (\pi_\R \times \pi_{TM}) \circ (\partial \times 0_{TM})
      = \Phi,
  \end{align*}
  using $\pi_\R \circ \partial = \id_\R$ and
  $\pi_{TM} \circ 0_{TM} = \id_{TM}$. Evaluating at $(0, v)$ and
  substituting (ii) gives $\Phi(0, v) = \pi_M(\Phi^{(1)}(0, v)) = \pi_M(v)$.
\end{proof}

\begin{lem}[Realization of variation fields]\label{lem:realization}
  Let $\gamma \in \Path(M; x, y)$, and let $W \colon \R \to TM$ be a smooth
  vector field along $\gamma$ (so $\pi_M \circ W = \gamma$) with
  $W(0) = 0_M(x)$ and $W(1) = 0_M(y)$. Then
  \begin{align*}
    \ad(P_W)(s, t) := \Phi\bigl(s, W(t)\bigr)
  \end{align*}
  defines a variation plot $P_W$ of $\gamma$
  (Definition~\ref{def:variation-plot}), and its variation field in the
  $s$-direction satisfies
  \begin{align*}
    \left. V_s \right|_{s=0} = [\ad(P_W), (0, t), \partial_s] = W.
  \end{align*}
\end{lem}

\begin{proof}
  The map $(s, t) \mapsto \Phi(s, W(t))$ is smooth, being a composite of the
  smooth maps $\Phi$ and $W$; by cartesian closedness of $\mathsf{Diff}$ it
  corresponds to a $1$-plot $P_W$ into the path space. By
  Lemma~\ref{lem:flow-initial}\,(i),
  $\ad(P_W)(0, t) = \Phi(0, W(t)) = \pi_M(W(t)) = \gamma(t)$, so $P_W(0) = \gamma$.
  Moreover the geodesic with initial velocity $0_M(x)$ is the constant path at
  $x$ (indeed $\dot\gamma \equiv 0$ satisfies
  $\nabla^{K}_\partial \dot\gamma = 0$ and $\dot\gamma(0) = 0_M(x)$, and by the
  uniqueness part of Theorem~\ref{thm:geodesic-existence} it is the only
  such geodesic); hence $\ad(P_W)(s, 0) = \Phi(s, 0_M(x)) = x$ for all $s$, and
  likewise $\ad(P_W)(s, 1) = y$. Thus $P_W$ is a variation plot in the sense of
  Definition~\ref{def:variation-plot}.

  Finally, for each fixed $t$ the curve
  $s \mapsto \ad(P_W)(s, t) = \Phi(s, W(t))$ is the geodesic with initial
  velocity $W(t)$ (Lemma~\ref{lem:flow-initial}), so, writing
  $\iota_t \colon \R \to \R^2$, $s \mapsto (s, t)$, for the slice map in the
  $s$-direction, Lemma~\ref{lem:flow-initial}\,(ii) gives
  \begin{align*}
   [\ad(P_W),(0,t),\partial_s]
     &= T(\ad(P_W))\circ T(\iota_t)\circ\partial\big|_0
      = T(\Phi)\circ(\partial\times 0_{TM})\,(0,\,W(t)) \\
     &= \Phi^{(1)}(0,\,W(t))
      = W(t).
  \end{align*}
  That is, $V_s|_{s=0} = W$.
\end{proof}

\begin{prop}[Converse of Proposition~\ref{prop:geodesic-critical-path}]%
\label{prop:critical-geodesic}
  If $\gamma \in \Path(M; x, y)$ is a critical path of the energy functional
  (Definition~\ref{def:critical-path}), then $\gamma$ is a geodesic on $M$;
  that is, $\nabla^{K}_\partial \dot\gamma = 0_M \circ \gamma$.
\end{prop}

\begin{proof}
  Set $Z := \nabla^{K}_\partial \dot\gamma = K \circ T(\dot\gamma) \circ \partial$,
  a vector field along $\gamma$ (Definition~\ref{def:covariant-derivative-along},
  Remark~\ref{rem:velocity-along}). It suffices to show $Z = 0_M \circ \gamma$.

  Let $W$ be any vector field along $\gamma$ vanishing at the endpoints. Taking
  the variation plot $P_W$ of Lemma~\ref{lem:realization}, we have
  $V_s|_{s=0} = W$; and since $P_W(0) = \gamma$,
  $V_t|_{s=0} = \dot\gamma$ (Remark~\ref{rem:variation-velocity}), whence
  $\nabla^{K}_{\partial_t} V_t|_{s=0} = \nabla^{K}_\partial \dot\gamma = Z$.
  The first variation formula (Lemma~\ref{lem:first-variation-formula}) and
  criticality of $\gamma$ give
  \begin{align}\label{eq:converse-star}
    \int_0^1 g(W, Z)\, dt
      = - \left. \frac{d(E \circ P_W)}{ds} \right|_{s=0}
      = 0
    \qquad \text{for every such } W.
  \end{align}
  The conclusion $Z \equiv 0$ now follows from the fundamental lemma of the
  calculus of variations, whose proof is identical to the manifold case once
  the test field is produced within the elastic setting: if $Z(t_0) \ne 0$ for
  some $t_0 \in (0,1)$, then $h := g(Z, Z)$ satisfies $h(t_0) > 0$ by
  definiteness, so $h > 0$ on an interval $(a,b) \ni t_0$; choosing
  $\varphi \in C^\infty(\R,\R)$ with $\operatorname{supp}\varphi \subset (a,b)$,
  $\varphi \ge 0$, $\varphi(t_0) > 0$ and setting $W := \varphi Z := \kappa_M
  \circ \langle \varphi, Z\rangle$ yields a vector field along $\gamma$
  vanishing at the endpoints (as $\varphi(0) = \varphi(1) = 0$), for which
  \eqref{eq:converse-star} and bilinearity give
  $0 = \int_0^1 \varphi\, h \, dt > 0$, a contradiction. Hence $Z(t) = 0_M(\gamma(t))$
  for all $t$, i.e.\ $\gamma$ is a geodesic.
\end{proof}

% ----------------------------------------
\subsection{Length-Minimizing Paths}\label{subsec:length}
% ----------------------------------------

\begin{defn}\label{def:length}
  A path $\gamma \in \Path(M; x, y)$ is \textbf{regular} if
  $g(\gamma)(t)(\partial, \partial) > 0$ for all $t$, and then its
  \textbf{length} is
  \begin{align*}
    L(\gamma)
      = \int_0^1 \bigl( g(\gamma)(t)(\partial, \partial) \bigr)^{1/2}\, dt,
  \end{align*}
  as in \cite{KSS}, in the notation of Notation~\ref{not:g-along-plot}.
\end{defn}

\begin{lem}\label{lem:energy-length}
  For every $\gamma \in \Path(M; x, y)$,
  \begin{align*}
    L(\gamma)^2 \le 2\, E(\gamma),
  \end{align*}
  with equality if and only if
  $t \mapsto g(\gamma)(t)(\partial, \partial)$ is constant.
\end{lem}

\begin{proof}
  Write $f(t) := (g(\gamma)(t)(\partial, \partial))^{1/2} \ge 0$.
  By the Cauchy--Schwarz inequality,
  \begin{align*}
    L(\gamma) = \int_0^1 f \cdot 1 \, dt
      \le \Bigl( \int_0^1 f^2 \, dt \Bigr)^{1/2}
          \Bigl( \int_0^1 1 \, dt \Bigr)^{1/2}
      = \bigl( 2 E(\gamma) \bigr)^{1/2}.
  \end{align*}
  Squaring gives the inequality; equality holds if and only if $f$ is
  proportional to the constant function $1$, i.e.\
  $f^2 = g(\gamma)(t)(\partial, \partial)$ is constant.
\end{proof}

\begin{lem}\label{lem:reparam}
  Let $\gamma \in \Path(M; x, y)$ be regular. Then there is a diffeomorphism
  $\psi \colon [0,1] \to [0,1]$ with $\psi(0) = 0$, $\psi(1) = 1$ such that
  $\widetilde\gamma := \gamma \circ \psi^{-1}$ satisfies
  \begin{align*}
    g(\widetilde\gamma)(t)(\partial, \partial) \equiv L(\gamma)^2,
    \qquad
    L(\widetilde\gamma) = L(\gamma),
    \qquad
    E(\widetilde\gamma) = \tfrac{1}{2} L(\gamma)^2.
  \end{align*}
\end{lem}

\begin{proof}
  Write $f(\tau) := g(\gamma)(\tau)(\partial, \partial)^{1/2} > 0$ and put
  $\psi(t) := L(\gamma)^{-1} \int_0^t f(\tau)\, d\tau$. Regularity gives
  $\psi' = L(\gamma)^{-1} f > 0$, so $\psi$ is a diffeomorphism of $[0,1]$ fixing
  the endpoints, with smooth inverse satisfying
  $(\psi^{-1})' = L(\gamma) / (f \circ \psi^{-1})$.

  Since $\psi^{-1}$ is a morphism of plots from $\widetilde\gamma$ to $\gamma$
  ($\gamma \circ \psi^{-1} = \widetilde\gamma$), functoriality of $T$ gives
  $T(\widetilde\gamma) = T(\gamma) \circ T(\psi^{-1})$, so by
  Remark~\ref{rem:tangent-structure-2}, for each $u$,
  \begin{align*}
    [\widetilde\gamma, u, \partial]
      = [\gamma,\ \psi^{-1}(u),\ (\psi^{-1})'(u)]
      = \kappa_M\bigl( (\psi^{-1})'(u),\ [\gamma, \psi^{-1}(u), \partial] \bigr),
  \end{align*}
  the second equality by
  $\kappa_M(r, [P, u, u']) = [P, u, r u']$
  (Definition~\ref{def:nat-trans}). Equivalently
  $\dot{\widetilde\gamma}
   = \kappa_M \circ \langle (\psi^{-1})',\ \dot\gamma \circ \psi^{-1} \rangle$;
  the only differentiation involved is that of $\psi^{-1}$ on the domain $\R$.

  By bilinearity of $g$,
  \begin{align*}
    g(\widetilde\gamma)(u)(\partial, \partial)
      = (\psi^{-1})'(u)^2 \, g(\gamma)(\psi^{-1}(u))(\partial, \partial)
      = \Bigl( \frac{L(\gamma)}{f(\psi^{-1}(u))} \Bigr)^{2} f(\psi^{-1}(u))^2
      = L(\gamma)^2 .
  \end{align*}
  Hence the speed is the constant $L(\gamma)$. Therefore
  $L(\widetilde\gamma) = \int_0^1 L(\gamma)\, du = L(\gamma)$ and
  $E(\widetilde\gamma) = \tfrac12 \int_0^1 L(\gamma)^2\, du = \tfrac12 L(\gamma)^2$.
\end{proof}

\begin{prop}\label{prop:minimizer-geodesic}
  Assume the standing hypotheses of \S\ref{subsec:energy}, so that
  Proposition~\ref{prop:critical-geodesic} holds. Let
  $\gamma \in \Path(M; x, y)$ be a regular path that minimizes length, i.e.\
  $L(\gamma) \le L(\sigma)$ for all $\sigma \in \Path(M; x, y)$. Then its
  constant-speed reparametrization $\widetilde\gamma$ (Lemma~\ref{lem:reparam})
  is a geodesic.
\end{prop}

\begin{proof}
  For any $\sigma \in \Path(M; x, y)$, Lemma~\ref{lem:energy-length} and
  minimality of $L(\gamma)$ give
  \begin{align*}
    E(\sigma) \ge \tfrac12 L(\sigma)^2 \ge \tfrac12 L(\gamma)^2
      = E(\widetilde\gamma),
  \end{align*}
  the last equality by Lemma~\ref{lem:reparam}. Hence $\widetilde\gamma$
  minimizes $E$ over $\Path(M; x, y)$. Consequently, for every variation plot
  $P$ of $\widetilde\gamma$ the smooth function $s \mapsto E(P(s))$ attains a
  minimum at $s = 0$ (as $P(s) \in \Path(M; x, y)$ for all $s$), so
  $\frac{d(E \circ P)}{ds}\big|_{s=0} = 0$; that is, $\widetilde\gamma$ is a
  critical path (Definition~\ref{def:critical-path}). By
  Proposition~\ref{prop:critical-geodesic}, $\widetilde\gamma$ is a geodesic.
\end{proof}

% ============================================================
\section{Example: Mapping Spaces}\label{sec:mapping-space}
% ============================================================

For elastic diffeological spaces, the explicit descriptions of geometric objects
in terms of domains that are standard in ordinary Riemannian geometry via local
coordinates are generally not available. The main reasons are that
$\colim_{P \in \D(M)}\hat{T}^2U_P \neq T^2M$ in general, and that vector fields
cannot be expressed as limits or colimits of vector fields on domains.
In this section we take mapping spaces as our example and examine the geometric
concepts constructed above in this concrete setting: the tangent structure,
connections and curvature, the Levi-Civita connection, and geodesics all lift
pointwise from the target.

% ----------------------------------------
\subsection{Tangent Structure on Mapping Spaces}\label{subsec:mapping-tangent}
% ----------------------------------------

The following results are due to Blohmann.

\begin{prop}[{\cite[Proposition~5.8]{B23}}]\label{prop:mapping-space-tangent-stable}
  Let $V$ be a tangent-stable elastic vector space.
  For any diffeological space $Z$, $C^\infty(Z, V)$ is elastic and
  \begin{align*}
    TC^\infty(Z, V) \cong C^\infty(Z, TV).
  \end{align*}
\end{prop}

\begin{prop}[{\cite[Corollary~5.9]{B23}}]\label{prop:mapping-space-manifolds}
  The mapping space $C^\infty(M, N)$ between manifolds is elastic, and
  \begin{align*}
    TC^\infty(M, N) \cong C^\infty(M, TN).
  \end{align*}
\end{prop}

In both cases the isomorphism is the canonical comparison map, which exists
for arbitrary $M$ and $N$: for $m \in M$, let
$\evl_m \colon C^\infty(M, N) \to N$ denote the evaluation, and define
\begin{align*}
  \Psi_{M, N} \colon TC^\infty(M, N) \longrightarrow C^\infty(M, TN),
  \qquad
  \Psi_{M, N}(v)(m) := T(\evl_m)(v),
\end{align*}
which is smooth by cartesian closedness of $\mathsf{Diff}$.
This motivates the following definition.

\begin{defn}\label{def:tangent-commuting}
A pair $(M, N)$ of diffeological spaces with $N$ elastic is called
\textbf{tangent-commuting} if $C^\infty(M, N)$ is elastic and the canonical map
$\Psi_{M, N} \colon TC^\infty(M, N) \to C^\infty(M, TN)$ is a diffeomorphism.
\end{defn}

Under this identification, $T(\evl_m)$ corresponds to the evaluation of
$C^\infty(M, TN)$ at $m$; iterating,
\begin{align*}
   T^k(\evl_m) \colon T^kC^\infty(M, N) \to T^kN
\end{align*}
corresponds to the evaluation of $C^\infty(M, T^kN)$ at $m$,
provided the pairs $(M, T^jN)$ for $j < k$ are tangent-commuting.

\begin{convention}\label{conv:hereditarily-tc}
  Throughout this section we assume that $(M, T^kN)$ is tangent-commuting for every $k \geq 0$ (where $T^0N = N$). When $N$ is a manifold this follows from Proposition~\ref{prop:mapping-space-manifolds}, since $T^kN$ is again a manifold; when $N$ is a tangent-stable elastic vector space it follows from Proposition~\ref{prop:mapping-space-tangent-stable}, since $TN \cong N \times N$ is again a tangent-stable elastic vector space.
  We write $\X := C^\infty(M, N)$ and suppress the identifications $T^k\X \cong C^\infty(M, T^kN)$ from the notation. Note also that, being a right adjoint by cartesian closedness, $C^\infty(M, -)$ preserves fibre products; in particular $T_2\X \cong C^\infty(M, T_2N)$ and $T\X \times_{\X} T\X \times_{\X} T\X \cong C^\infty(M, TN \times_N TN \times_N TN)$.
\end{convention}

This assumption is used repeatedly to identify
$T^kC^\infty(M,N)$
with
$C^\infty(M,T^kN)$,
thereby allowing all higher-order constructions (connections, curvature and geodesic sprays) to be evaluated pointwise.

Since the tangent structure $(\pi_N, 0_N, \kappa_N, +_N, \tau_N, \lambda_N)$ on $N$ consists of natural transformations, and post-composition $C^\infty(M, -)$ is a covariant functor, the tangent structure on $\X$ is obtained by applying $C^\infty(M, -)$ to that of $N$. Explicitly:
\begin{description}
  \item[Projection]
    $\pi_{\X}(v) = \pi_N \circ v$,\quad
    $v \in C^\infty(M, TN)$.
  \item[Zero section]
    $0_{\X}(f) = 0_N \circ f$,\quad
    $f \in C^\infty(M, N)$.
  \item[Addition]
    $(v_1 +_{\X} v_2)(x) = v_1(x) +_N v_2(x)$,\quad
    $x \in M$.
  \item[Scalar multiplication]
    $\kappa_{\X}(r, v)(x) = \kappa_N(r, v(x))$,\quad
    $r \in \R$.
  \item[Symmetric structure]
    $\tau_{\X}(u) = \tau_N \circ u$,\quad
    $u \in C^\infty(M, T^2N)$.
  \item[Vertical lift]
    $\lambda_{\X}(u) = \lambda_N \circ u$.
\end{description}
The axioms of a tangent structure are equations between composites of these
natural transformations, and are therefore preserved by the functor
$C^\infty(M, -)$.

% ----------------------------------------
\subsection{Connections and Curvature on Mapping Spaces}%
\label{subsec:mapping-connection}
% ----------------------------------------

Connections on mapping spaces likewise lift from the target space.

\begin{prop}\label{prop:mapping-space-connection}
 Let $K_N$ and $H_N$ be a vertical connection and a horizontal connection on
 $N$, respectively. Define
 $K_{\X} \colon T^2\X \to T\X$ and
 $H_{\X} \colon T\X \times_{\X} T\X \to T^2\X$ by
 \begin{align*}
   K_{\X}(u) &:= K_N \circ u,
     \qquad u \in C^\infty(M, T^2N),\\
   H_{\X}(v, w) &:= H_N \circ \langle v, w \rangle,
     \qquad \langle v, w \rangle \in C^\infty(M, TN \times_N TN).
 \end{align*}
 Then $K_{\X}$ is a vertical connection on $\X$ and $H_{\X}$ is a horizontal
 connection on $\X$. If $K_N$ is torsion-free, so is $K_{\X}$; if $(K_N, H_N)$
 is a full connection, so is $(K_{\X}, H_{\X})$.
\end{prop}

\begin{proof}
  Under the identifications of Convention~\ref{conv:hereditarily-tc}, $K_{\X} = C^\infty(M, K_N)$, $H_{\X} = C^\infty(M, H_N)$, and every structure map appearing in Definitions~\ref{def:vertical-connection}, \ref{def:horizontal-connection} and~\ref{def:full-connection} is the image under $C^\infty(M, -)$ of the corresponding map for $N$. 
  Each axiom, as well as torsion-freeness, is an equation between composites of such maps,
  and is therefore inherited from $N$ by functoriality; the scalar axioms, which involve the parameter $\R$, are likewise verified pointwise in $m \in M$.
\end{proof}

\begin{prop}\label{prop:mapping-space-effective}
   If $K_N$ is an effective vertical connection on $N$ (Definition~\ref{def:effective-vertical-connection}), then $K_{\X}$ is also effective.
\end{prop}

\begin{proof}
  Under the same identifications (using that $C^\infty(M, -)$ preserves
  fibre products),
  \begin{align*}
    \Theta_{\X} = \langle \pi_{T\X}, T(\pi_{\X}), K_{\X} \rangle
      = C^\infty(M, \Theta_N);
  \end{align*}
  since functors preserve isomorphisms,
  $C^\infty(M, \Theta_N^{-1})$ is a smooth inverse of $\Theta_{\X}$, so
  $K_{\X}$ is effective.
\end{proof}

The curvature (Section~\ref{sec:curvature}) also localizes.

\begin{prop}\label{prop:mapping-space-curvature}
  Let $K_N$ be a vertical connection on $N$. Under the identification $T^3\X \cong C^\infty(M, T^3N)$, the curvature morphism of $K_{\X}$ (Definition~\ref{def:curvature-morphism}) is given by post-composition with that of $K_N$:
  \begin{align*}
    C_{K_{\X}}(u) = C_{K_N} \circ u,
    \qquad u \in C^\infty(M, T^3N).
  \end{align*}
  In particular, $K_{\X}$ is flat if and only if $K_N$ is flat.
\end{prop}

\begin{proof}
  Both composites $K_N \circ T(K_N) \circ \tau_{TN}$ and $K_N \circ T(K_N)$, as well as the fibrewise subtraction in the abelian group structure of $\pi_N \colon TN \to N$, are built from structure maps of $N$ that lift by post-composition; in particular the fibrewise addition and inverse of $\pi_{\X}$ are obtained by applying
  $C^\infty(M, -)$ to those of $\pi_N$. 
  Hence $C_{K_{\X}} = C^\infty(M, C_{K_N})$ by the same argument as in Proposition~\ref{prop:mapping-space-connection}.
  If $C_{K_N} = 0$, then $C_{K_{\X}}(u) = C_{K_N} \circ u = 0$ for every $u$.
  Conversely, if $C_{K_{\X}} = 0$, then evaluating at constant maps $u \equiv w$ for $w \in T^3N$ gives $C_{K_N}(w) = 0$.
\end{proof}

\begin{ex}\label{ex:flat-mapping-space}
  For $N = \R^n$ with the canonical vertical connection $K_0$ of Example~\ref{ex:euclidean-vertical-connection}, $C_{K_0} = 0$ by Example~\ref{ex:flat-euclidean}. Hence, for any manifold $M$, the connection $K_{C^\infty(M, \R^n)}$ on the mapping space $C^\infty(M, \R^n)$ is flat.
\end{ex}

% ----------------------------------------
\subsection{The Levi-Civita Connection on Mapping Spaces}%
\label{subsec:mapping-lc}
% ----------------------------------------

For the remainder of this section, $M$ is a closed manifold equipped with a
volume form $\vol_M$, and
$(N, g_N)$ is an elastic Riemannian diffeological space.

\begin{defn}\label{def:g-map}
  The \textbf{mapping-space metric} $g_{\mathit{map}}$ on $\X = C^\infty(M, N)$ is defined, for $v, w \in T\X$ lying in a common
  fibre, by
  \begin{align*}
    g_{\mathit{map}}(v, w)
      := \int_M g_N\bigl(T(\evl_m)(v),\, T(\evl_m)(w)\bigr)\, \vol_M(m).
  \end{align*}
  In the plot notation of Notation~\ref{not:g-along-plot} this reads $g_{\mathit{map}}(P)(u)(v, w) = \int_M g_N(\evl_m \circ P)(u)(v, w)\,\vol_M(m)$, as in \cite{KSS}.
\end{defn}

\begin{lem}\label{lem:g-map-metric}
  $g_{\mathit{map}}$ is a Riemannian metric on $\X$ in the sense of
  Definition~\ref{def:metric}; in particular it is definite.
\end{lem}

\begin{proof}
  Write $\mathcal{I} \colon C^\infty(M, \R) \to \R$ for the integration against $\vol_M$, a smooth linear map, and $\mathcal{G} \colon T_2\X \to C^\infty(M, \R)$ for the pointwise metric evaluation, $\mathcal{G}(v, w)(m) := g_N(T(\evl_m)(v), T(\evl_m)(w))$, which is smooth by cartesian closedness. 
  Then $g_{\mathit{map}} = \mathcal{I} \circ \mathcal{G}$ is smooth. 
  Bilinearity and symmetry hold pointwise in $m$ and are preserved by the linear map $\mathcal{I}$; positivity follows since the integrand is non-negative. For definiteness, suppose $g_{\mathit{map}}(v, v) = 0$. 
  The function $m \mapsto g_N(T(\evl_m)(v), T(\evl_m)(v))$ is smooth, non-negative, and has vanishing integral over the closed manifold $M$, hence vanishes identically; definiteness of $g_N$ then gives $T(\evl_m)(v) = 0$ for every $m$, i.e.\ $\Psi_{M,N}(v) = 0$, and so $v = 0$.
\end{proof}

Let $K_N$ be a vertical Levi-Civita connection for $g_N$.

\begin{prop}\label{prop:mapping-space-levi-civita}
   $K_{\X}$ is a vertical Levi-Civita connection for $g_{\mathit{map}}$.
   By Theorem~\ref{thm:uniqueness-cov-deriv}, its covariant derivative $\nabla^{K_{\X}}$ is the unique Levi-Civita covariant derivative for $g_{\mathit{map}}$.
\end{prop}

\begin{proof}
   Torsion-freeness of $K_{\X}$ was shown in Proposition~\ref{prop:mapping-space-connection}. We verify the metric
   compatibility condition \eqref{eq:metric-compatibility} for $g_{\mathit{map}}$.

   For $m \in M$ write
   $h_m := g_N \circ (T(\evl_m) \times_{\evl_m} T(\evl_m))
   \colon T_2\X \to \R$, so that
   $g_{\mathit{map}} = \mathcal{I} \circ \mathcal{G}$ with
   $\mathcal{G}(\xi)(m) = h_m(\xi)$ as in Lemma~\ref{lem:g-map-metric}.
   Since $\mathcal{I}$ is smooth and linear, under the identification
   $TC^\infty(M, \R) \cong C^\infty(M, T\R) \cong C^\infty(M, \R \times \R)$
   its tangent map acts componentwise, whence
   $d(\mathcal{I} \circ \mathcal{G})
    = \mathcal{I} \circ \bigl(m \mapsto d(h_m)\bigr)$; that is,
   \begin{align}\label{eq:diff-under-integral}
     d(g_{\mathit{map}})(W)
       = \int_M d(h_m)(W)\, \vol_M(m),
     \qquad W \in T(T_2\X).
   \end{align}
   Next, $h_m = g_N \circ T_2(\evl_m)$ with
   $T_2(\evl_m) := T(\evl_m) \times T(\evl_m) \colon T_2\X \to T_2N$, so the
   chain rule $d(f \circ h) = d(f) \circ T(h)$ (immediate from
   $d = \pr_2 \circ T$ and functoriality) gives
   $d(h_m) = d(g_N) \circ T(T_2(\evl_m))$, where $T(T_2(\evl_m))$ is
   identified with $T^2(\evl_m) \times T^2(\evl_m)$ on
   $T(T_2\X) \cong T^2\X \times_{T\X} T^2\X$. Substituting the metric
   compatibility condition \eqref{eq:metric-compatibility} for $(N, g_N, K_N)$,
   \begin{align*}
     d(h_m)(W_1, W_2)
       = g_N\bigl(K_N T^2(\evl_m) W_1,\; \pi_{TN} T^2(\evl_m) W_2\bigr)
       + g_N\bigl(\pi_{TN} T^2(\evl_m) W_1,\; K_N T^2(\evl_m) W_2\bigr).
   \end{align*}
   By Convention~\ref{conv:hereditarily-tc}, $T^2(\evl_m)$ is the evaluation
   of $C^\infty(M, T^2N)$ at $m$, so
   $K_N \circ T^2(\evl_m) = T(\evl_m) \circ K_{\X}$ and
   $\pi_{TN} \circ T^2(\evl_m) = T(\evl_m) \circ \pi_{T\X}$. Hence
   \begin{align*}
     d(h_m)(W_1, W_2)
       = h_m\bigl(K_{\X}(W_1),\, \pi_{T\X}(W_2)\bigr)
       + h_m\bigl(\pi_{T\X}(W_1),\, K_{\X}(W_2)\bigr),
   \end{align*}
   and integrating over $M$ via \eqref{eq:diff-under-integral},
   \begin{align*}
     d(g_{\mathit{map}})
       = +_{\R} \circ
         \langle g_{\mathit{map}} \circ (K_{\X} \times_{\pi} \pi_{T\X}),\;
                g_{\mathit{map}} \circ (\pi_{T\X} \times_{\pi} K_{\X}) \rangle,
   \end{align*}
   which is \eqref{eq:metric-compatibility} on $\X$.
\end{proof}

% ----------------------------------------
\subsection{Geodesics on Mapping Spaces}\label{subsec:mapping-geodesics}
% ----------------------------------------

We keep the hypotheses of \S\ref{subsec:mapping-lc}. For a path
$\gamma \colon \R \to \X$ and $m \in M$, write
$\gamma_m := \evl_m \circ \gamma \colon \R \to N$, a path on $N$.

\begin{prop}\label{prop:mapping-space-geodesic}
  A path $\gamma \colon \R \to \X$ is a geodesic for $K_{\X}$ if and only if
  $\gamma_m$ is a geodesic on $N$ for $K_N$ for every $m \in M$.
\end{prop}

\begin{proof}
  By naturality of the identifications, for every $m \in M$,
  \begin{align*}
    T(\evl_m)\bigl(\nabla^{K_{\X}}_\partial \dot\gamma\bigr)
      &= T(\evl_m) \circ K_{\X} \circ T(\dot\gamma) \circ \partial
       = K_N \circ T^2(\evl_m) \circ T(\dot\gamma) \circ \partial \\
      &= K_N \circ T\bigl(T(\evl_m) \circ \dot\gamma\bigr) \circ \partial
       = K_N \circ T(\dot\gamma_m) \circ \partial
       = \nabla^{K_N}_\partial \dot\gamma_m,
  \end{align*}
  using $T(\evl_m) \circ \dot\gamma
  = T(\evl_m) \circ T(\gamma) \circ \partial = T(\gamma_m) \circ \partial
  = \dot\gamma_m$, and likewise
  $T(\evl_m) \circ 0_{\X} \circ \gamma = 0_N \circ \gamma_m$.
  Hence, if $\gamma$ is a geodesic, applying $T(\evl_m)$ to
  $\nabla^{K_{\X}}_\partial \dot\gamma = 0_{\X} \circ \gamma$ shows that each
  $\gamma_m$ is a geodesic. Conversely, if every $\gamma_m$ is a geodesic,
  then $T(\evl_m) \circ \nabla^{K_{\X}}_\partial \dot\gamma
  = T(\evl_m) \circ 0_{\X} \circ \gamma$ for every $m$; since the family
  $(T(\evl_m))_{m \in M}$ corresponds under $\Psi_{M,N}$ to the evaluations of
  $C^\infty(M, TN)$, it is jointly injective, so
  $\nabla^{K_{\X}}_\partial \dot\gamma = 0_{\X} \circ \gamma$.
\end{proof}

The geodesic flow, and hence the existence theory of
\S\ref{subsec:geodesic-existence}, also transfers pointwise.

\begin{prop}\label{prop:mapping-space-flow}
  Suppose $N \in \mathsf{Elast}_{\mathrm{curve}}$ and let $(K_N, H_N)$ be a full
  connection on $N$ whose geodesic spray $S_N$ is complete, with geodesic flow
  $\Phi_N \colon \R \times TN \to N$. Then the geodesic spray $S_{\X}$ of
  $(K_{\X}, H_{\X})$ is complete, with geodesic flow
  \begin{align*}
    \Phi_{\X} \colon \R \times T\X \to \X,
    \qquad
    \Phi_{\X}(s, v)(m) := \Phi_N\bigl(s,\, T(\evl_m)(v)\bigr).
  \end{align*}
  Moreover, for every $v_0 \in T\X$ there is a unique geodesic
  $\gamma \colon \R \to \X$ with $\dot\gamma(0) = v_0$, namely
  $\gamma = \Phi_{\X} \circ \iota_{v_0}$.
\end{prop}

\begin{proof}
  The map $(s, v, m) \mapsto \Phi_N(s, T(\evl_m)(v))$ is smooth, so
  $\Phi_{\X}$ is smooth by cartesian closedness. Under the identifications,
  $S_{\X} = C^\infty(M, S_N)$ and
  $\evl_m \circ \Phi_{\X} = \Phi_N \circ (\id_\R \times T(\evl_m))$;
  applying $T$ and $T^2$ to the latter and using naturality of $0$, one
  obtains
  \begin{align*}
    T(\evl_m) \circ \Phi_{\X}^{(1)}
      = \Phi_N^{(1)} \circ (\id_\R \times T(\evl_m)),
    \qquad
    T^2(\evl_m) \circ \Phi_{\X}^{(2)}
      = \Phi_N^{(2)} \circ (\id_\R \times T(\evl_m)).
  \end{align*}
  The initial and differential conditions for $\Phi_{\X}$
  (Definition~\ref{def:nth-solution}) then follow from those for $\Phi_N$ by
  the joint injectivity of the families $(T(\evl_m))_m$ and
  $(T^2(\evl_m))_m$, as in Proposition~\ref{prop:mapping-space-geodesic}.
  Hence $S_{\X}$ is complete with geodesic flow $\Phi_{\X}$.

  For the final claim, existence: as in
  Theorem~\ref{thm:geodesic-existence}, the slice
  $\gamma = \Phi_{\X} \circ \iota_{v_0}$ satisfies
  $\dot\gamma(0) = \Phi_{\X}^{(1)}(0, v_0) = v_0$, and each
  $\gamma_m = \Phi_N(\cdot, T(\evl_m)(v_0))$ is a geodesic on $N$
  (Lemma~\ref{lem:flow-initial}), so $\gamma$ is a geodesic by
  Proposition~\ref{prop:mapping-space-geodesic}. Uniqueness: if $\gamma$ and
  $\sigma$ are geodesics on $\X$ with
  $\dot\gamma(0) = \dot\sigma(0) = v_0$, then for each $m$ the paths
  $\gamma_m$ and $\sigma_m$ are geodesics on $N$ with the same initial
  velocity $T(\evl_m)(v_0)$, so $\gamma_m = \sigma_m$ by
  Theorem~\ref{thm:geodesic-existence} applied to $N$; since the
  evaluations $(\evl_m)_m$ are jointly injective, $\gamma = \sigma$.
\end{proof}

\begin{rem}\label{rem:mapping-space-lc-membership}
  Proposition~\ref{prop:mapping-space-flow} provides existence and uniqueness
  of geodesics on $\X$ without settling whether
  $\X \in \mathsf{Elast}_{\mathrm{curve}}$, i.e.\ whether $\R$ is a curve object
  for $\X$: the uniqueness above is derived pointwise from
  $N \in \mathsf{Elast}_{\mathrm{curve}}$ rather than from a curve-object
  property of $\X$. Whether $N \in \mathsf{Elast}_{\mathrm{curve}}$ implies
  $C^\infty(M, N) \in \mathsf{Elast}_{\mathrm{curve}}$ is left open.
\end{rem}

Collecting the preceding propositions, we obtain the following theorem.

\begin{thm}\label{thm:summary-of-mapping-space}
Let $M$ be a closed manifold equipped with a volume form $\vol_M$ and let $(N,g_N)$ be an elastic Riemannian diffeological space such that $(M, T^kN)$ is tangent-commuting for every $k \ge 0$. Write $\X = C^\infty(M,N)$.
\begin{enumerate}
  \item $\X$ is elastic and its tangent structure is obtained pointwise from
    that of $N$.
  \item $g_{\mathit{map}}$ is a Riemannian metric on $\X$.
  \item If $K_N$ is a vertical connection on $N$, then $K_\X$ is a vertical connection on $\X$; $K_\X$ is torsion-free (resp.\ flat, effective) whenever $K_N$ is, and its curvature is $C_{K_N}\circ(-)$.
  \item If $K_N$ is a Levi-Civita connection for $g_N$, then $K_\X$ is a Levi-Civita connection for $g_{\mathit{map}}$.
  \item If $(K_N,H_N)$ is a full connection whose geodesic spray is complete, then the geodesic spray of $(K_\X,H_\X)$ is complete. If moreover $N \in \mathsf{Elast}_{\mathrm{curve}}$, every $v_0 \in T\X$ determines a unique geodesic.
\end{enumerate}
\end{thm}

% ============================================================
\section{Perspectives}\label{sec:perspectives}
% ============================================================

The present work establishes a first framework for Riemannian geometry on
elastic diffeological spaces. Several natural directions arise from this
construction, including affine geometry on elastic spaces, Lie-theoretic
applications, and categorical formulations of Riemannian geometry; we
briefly describe some of them below. A detailed treatment of the first two
will appear in a separate paper.

% ------------------------------------------------------------
\subsection{Connections on elastic vector spaces}\label{subsec:elastic-vs}
% ------------------------------------------------------------

A natural next step is a systematic theory of affine connections on
tangent-stable elastic vector spaces, for which the Euclidean computations
of Section~\ref{sec:connections} serve as the finite-dimensional model. One
expects vertical connections on such spaces to be classified by
Christoffel-type operators --- bilinear correction terms added to a
canonical connection --- and, in particular, to be automatically effective,
so that vertical and full connections coincide there. A further step is
connections on diffeological manifolds modelled on such spaces; this raises
the questions of whether elasticity is a local property and of how to glue
local Christoffel data in the absence of smooth partitions of unity.

% ------------------------------------------------------------
\subsection{Elastic groups}\label{subsec:elastic-group}
% ------------------------------------------------------------

Another promising direction is the study of left-invariant Riemannian
metrics on elastic groups. For left-invariant vector fields the derivative
terms in the Koszul formula vanish, and one expects the geodesic equation
to reduce to an Euler--Arnold equation formulated using only the metric and
the tangent-category bracket --- avoiding the dual of the Lie algebra, the
coadjoint action, and the inertia operator that are the usual analytic
obstacles in infinite dimensions. Loop groups, which are elastic with
pointwise geodesics by the results of Section~\ref{sec:mapping-space}, and
their central extensions, such as affine Kac--Moody groups and the Virasoro
group, provide natural test cases, with geometric interpretations of
integrable equations such as KdV and Camassa--Holm as a long-term aim. Such
a theory would extend classical Lie group geometry to elastic spaces and
clarify the interaction between algebraic and geometric structures in this
setting.

% ------------------------------------------------------------
\subsection{Tangent-categorical formulation}\label{subsec:tangent-core}
% ------------------------------------------------------------

Since many constructions in this paper rely only on the tangent-categorical
structure, it is natural to ask over which tangent categories the theory
survives. The metric compatibility condition, the Koszul formula, and the
uniqueness of the Levi-Civita covariant derivative use only the tangent
structure, the Lie bracket, and the ring operations, and should carry over
to Cartesian tangent categories with scalar multiplication over a
commutative ring object; by contrast, positivity, the musical maps, and the
geodesic theory require additional structure --- an ordered ring, a
cotangent identification, and a curve object, respectively. Isolating the
precise axioms needed for each of these layers would reveal the categorical
essence of elastic Riemannian geometry and facilitate comparisons with
existing axiomatic frameworks.

% ------------------------------------------------------------
\subsection{Further geometric problems}\label{subsec:further-problems}
% ------------------------------------------------------------

Several foundational questions remain open, beginning with those of Open
Question~\ref{prob:open-questions}: the existence of a Levi-Civita
connection for a strongly non-degenerate metric, and the uniqueness of the
connection itself rather than of its covariant derivative --- the
obstruction being that the covariant derivative need not determine the
connection in the absence of local triviality. On mapping spaces, it is
open whether the curve-object property passes from the target $N$ to
$C^\infty(M, N)$. For global geometry, natural next goals are a Gauss-type
lemma showing that geodesics are locally length-minimizing, analogues of
the Hopf--Rinow theorem, and comparison results based on curvature
invariants.

\medskip
The directions outlined above indicate that elastic Riemannian geometry has
the potential to develop into a broad framework encompassing differential
geometry, Lie theory, and categorical geometry. These topics will be
investigated in future work.

% ============================================================
\appendix
\section{Tangent categories}\label{sec:tangent-categories}

In this appendix, we summarize the basic categorical concepts regarding tangent categories used in the main text.
For detailed introductions and proofs, we refer the reader to \cite{AB, CC14, CC15, M} and the references therein.
Note that in this appendix, following convention, the objects of an abstract tangent category $\mathcal{C}$ are denoted by $X, Y, Z$.

\subsection{Bundles with algebraic structure and symmetric structures}\label{sec:first-definitions}

Consider a category $\mathcal{C}$ that has a terminal object $*$ and such that for any object $X$, the overcategory $\mathcal{C}\downarrow X$ has finite products (this implies that $\mathcal{C} \cong \mathcal{C} \downarrow *$ itself also has finite products).

\begin{defn}\label{def:bundle-abelian-groups}
  A bundle of abelian groups $p\colon A \to X$ is an abelian group object in $\mathcal{C}\downarrow X$.
\end{defn}

This data consists of the bundle projection $p\colon A \to X$, along with morphisms for addition $+\colon A \times_X A \to A$, the zero-section $0\colon X \to A$, and the inverse $\iota\colon A \to A$, satisfying the usual axioms of an abelian group (associativity, identity, commutativity, etc.). No local triviality is assumed.

\begin{defn}\label{def:bundle-morphism}
  Let $p\colon A \to X$ and $p'\colon A' \to X'$ be bundles of abelian groups with addition $+$ and $+'$, respectively.
  A morphism of bundles is a commutative diagram in the category $\mathcal{C}$ of the following form:
  $$
    \begin{tikzcd}
      A \ar[r, "\phi"] \ar[d, "p"'] & A' \ar[d, "p'"] \\
      X \ar[r, "f"'] & X'
    \end{tikzcd}
  $$
  Furthermore, if the following diagram commutes, $\phi$ is called a morphism of bundles of abelian groups:
  $$
    \begin{tikzcd}
      A \times_X A \ar[r, "\phi \times_f \phi"] \ar[d, "+"'] & A' \times_{X'} A' \ar[d, "+'"] \\
      A \ar[r, "\phi"'] & A'
    \end{tikzcd}
  $$
\end{defn}

For a commutative ring object $R$ in $\mathcal{C}$, the projection $X \times R \to X$ becomes a ring object in $\mathcal{C} \downarrow X$.
A pair consisting of a bundle of abelian groups $p\colon A \to X$ and a scalar multiplication $\kappa\colon R \times A \to A$ satisfying the axioms of a left module is called a bundle of $R$-modules.

\begin{defn}\label{def:symmetric-structure}
  Consider a covariant endofunctor $T\colon \mathcal{C} \to \mathcal{C}$ and a natural transformation $\tau\colon T^2 \to T^2$.
  Let $\tau_{12} \coloneqq \tau T$ and $\tau_{23} \coloneqq T \tau$ be the trivial extensions to natural transformations $T^3 \to T^3$.
  If $\tau \circ \tau = 1$ and the braid relation
  $\tau_{12} \circ \tau_{23} \circ \tau_{12} = \tau_{23} \circ \tau_{12} \circ \tau_{23}$
  is satisfied, then $\tau$ is called a symmetric structure.
\end{defn}

\begin{defn}\label{def:preserves-fiber-products}
  For a bundle of abelian groups $p\colon A \to X$, the covariant endofunctor $T\colon \mathcal{C} \to \mathcal{C}$ is said to preserve fibre products of $p\colon A \to X$ if the natural morphisms
  $$
    \nu_{k,X} \colon T(A \times_X \dots \times_X A)
      \to TA \times_{TX} \dots \times_{TX} TA
  $$
  are all isomorphisms for any $k \geq 1$.
\end{defn}

\subsection{Axioms of tangent categories}\label{sec:RosickysAxioms}

\begin{defn}\label{def:tangent-structure-axioms}
  A tangent structure on a category $\mathcal{C}$ consists of a covariant endofunctor $T\colon \mathcal{C} \to \mathcal{C}$ and the following natural transformations:
  \makeatletter
  $$
  \begin{alignedat}{2}
    & \text{projection}\quad &  \pi &\colon T \to 1 \\
    &\text{zero-section}\quad & 0 &\colon 1 \to T \\
    & \text{addition}\quad & +&\colon T_2 \to T \\
    &\text{vertical lift}\quad & \lambda&\colon T \to T^2 \\
    &\text{symmetric structure}\quad & \tau &\colon T^2 \to T^2.
  \end{alignedat}
  $$
  \makeatother
  The axioms they must satisfy are as follows:
  \begin{itemize}
    \item The projection $\pi\colon T \to 1$ is a bundle of abelian groups over $1$ with zero-section $0$ and addition $+$.
    \item All pullbacks $T_k \coloneqq T \times_1 \dots \times_1 T$ exist and are preserved pointwise by $T$.
    \item The natural transformation $\tau$ is a symmetric structure and a morphism of bundles of abelian groups.
    \item The following diagrams commute, and in particular, the left diagram is a morphism of bundles of abelian groups:
      $$
        \begin{tikzcd}
          T  \ar[r, "\lambda"] \ar[d, "\pi"'] & T^2 \ar[d, "\pi T"] \\
          1  \ar[r, "0"'] & T
        \end{tikzcd}
        \qquad
        \begin{tikzcd}
          T \ar[r, "\lambda"] \ar[d, "\lambda"'] & T^2   \ar[d, "\lambda T"] \\
          T^2  \ar[r, "T\lambda"'] & T^3
        \end{tikzcd}
      $$
    \item The following diagrams commute with respect to the vertical lift and the symmetric structure:
      $$
        \begin{tikzcd}
          & T \ar[dl, "\lambda"'] \ar[dr, "\lambda"] & \\
          T^2  \ar[rr, "\tau"'] & & T^2
        \end{tikzcd}
        \qquad
        \begin{tikzcd}
          T^2  \ar[r, "T\lambda"] \ar[d, "\tau"'] & T^3  \ar[r, "\tau T"] & T^3 \ar[d, "T\tau"] \\
          T^2  \ar[r, "\lambda T"'] & T^3 \ar[r, "\tau"'] & T^3
        \end{tikzcd}
      $$
    \item The diagram formed by the following natural transformation $\lambda_2$ is a pointwise pullback:
      $$
        \begin{tikzcd}
          T_2  \ar[r, "T0 \times_0 \lambda"] \ar[d, "\pi \circ \pr_1"']
            & T_2T \ar[r, "+_T"]
            & T^2 \ar[r, "\tau"]
            & T^2 \ar[d, "T\pi"] \\
          1  \ar[rrr, "0"'] & & & T
        \end{tikzcd}
      $$
  \end{itemize}
\end{defn}

\subsection{Cartesian tangent categories and scalar multiplication}\label{sec:CatTanScal}

Suppose the tangent category $\mathcal{C}$ has finite products.

\begin{defn}\label{def:cartesian-tangent-category}
  A tangent category $(\mathcal{C}, T)$ is called a Cartesian tangent category if $T* \cong *$ and the natural morphisms
  $$
    \chi_{X,Y}\colon T(X \times Y) \longrightarrow TX \times TY
  $$
  are isomorphisms for any objects $X, Y$.
\end{defn}

In a Cartesian tangent category, for a morphism $f\colon X \times Y \to Z$, the partial tangent morphisms with respect to each variable can be defined as follows:
$$
\begin{aligned}
  T_{(1)} f \colon TX \times Y
    &\xrightarrow{\id_{TX} \times 0_Y} TX \times TY
     \xrightarrow{\chi_{X,Y}^{-1}} T(X \times Y)
     \xrightarrow{Tf} TZ, \\
  T_{(2)} f \colon X \times TY
    &\xrightarrow{0_X \times \id_{TY}} TX \times TY
     \xrightarrow{\chi_{X,Y}^{-1}} T(X \times Y)
     \xrightarrow{Tf} TZ.
\end{aligned}
$$

\begin{prop}[{\cite[Proposition 2.10]{CC14}}]\label{prop:decomposition-tangent-morphism}
  For a morphism $f\colon X \times Y \to Z$ in a Cartesian tangent category $(\mathcal{C}, T)$, the following holds:
  $$
  \begin{aligned}
    Tf = T_{(1)} f \circ (\id_{TX} \times \pi_Y) \circ \chi_{X,Y}
         +_{TZ} T_{(2)} f \circ (\pi_X \times \id_Y) \circ \chi_{X,Y}.
  \end{aligned}
  $$
\end{prop}

Fix a commutative ring object $R$ in a Cartesian tangent category $\mathcal{C}$.
The natural transformation $1 \times R \to 1$ inherits the structure of a ring object in the category $\operatorname{End}(\mathcal{C}) \downarrow 1$ from $R$.
An $R$-module structure on $T$ is a $(1 \times R \to 1)$-module structure on $\pi$.
This is given by a natural morphism $\kappa\colon R \times T \to T$.
Furthermore, if the natural morphism $TR \to R \times R$ is an isomorphism of $R$-modules, $R$ is called tangent-stable.

\begin{defn}\label{def:scalar-multiplication}
  A Cartesian tangent category with scalar $R$-multiplication is a tuple consisting of a Cartesian tangent category $\mathcal{C}$, a commutative ring object $R$, and a tangent-stable $R$-module structure $\kappa$ on $T$ that makes the following diagrams commute:
  $$
    \begin{tikzcd}[column sep=large]
      R \times TX \ar[r, "\id_R \times \lambda_X"] \ar[d, "\kappa_X"']
        & R \times T^2 X \ar[d, "\kappa_{TX}"] \\
      TX \ar[r, "\lambda_X"'] & T^2 X
    \end{tikzcd}
  $$
  $$
    \begin{tikzcd}[column sep=6em]
      TR \times TX \ar[r, "T_{(1)}\kappa_X"] \ar[d, "{(\pi_R, \mathsf{Val}_R) \times \id_{TX}}"', "\cong"]
        &[4em] T^2X \\
      R \times R \times TX
        \ar[r, "{(\kappa_X \circ (\pr_1,\pr_3),\, \kappa_X \circ (\pr_2,\pr_3))}"']
        & T_2 X \ar[u, "\lambda_{2,X}"']
    \end{tikzcd}
  $$
  $$
    \begin{tikzcd}[column sep=3em]
      R \times T^2 X \ar[r, "T_{(2)}\kappa_X"] \ar[d, "\id_R \times \tau_X"'] & T^2 X \\
      R \times T^2 X \ar[r, "\kappa_{TX}"'] & T^2 X \ar[u, "\tau_X"']
    \end{tikzcd}
  $$
  Here, $\mathsf{Val}_R \colon TR \to R$ is the projection onto the second component of $TR \cong R \times R$.
\end{defn}

\subsection{Differential bundles and their morphisms}

The concept of a \textbf{differential bundle} in a tangent category was introduced in \cite{CC18}.
A differential bundle over an object $M$ is a tuple $\mathsf{q} = (q,\, z_q,\, \lambda_q)$ consisting of a morphism $q \colon E \to M$ (projection), $z_q \colon M \to E$ (zero-section), and a lift $\lambda_q \colon E \to TE$, satisfying various conditions compatible with the axioms of a tangent category (see \cite{CC18} for details).
The tangent bundle $\pi_M \colon TM \to M$ is a prototypical example of a differential bundle.

\begin{defn}[\cite{CC18}]\label{def:differential-bundle-morphism}
Let $\mathsf{q} = (q,\, z_q,\, \lambda_q)$ and $\mathsf{q}' = (q',\, z_{q'},\, \lambda_{q'})$ be differential bundles over a tangent category $\mathcal{C}$.
\begin{enumerate}
  \item A \textbf{bundle morphism} between them is a pair of morphisms $f_1 \colon E \to E'$ and $f_0 \colon M \to M'$ satisfying $q' \circ f_1 = f_0 \circ q$:
    $$
      \begin{tikzcd}
        E \ar[r, "f_1"] \ar[d, "q"'] & E' \ar[d, "q'"] \\
        M \ar[r, "f_0"'] & M'
      \end{tikzcd}
    $$
  \item A bundle morphism $(f_1, f_0)$ is said to be \textbf{linear} if it additionally satisfies $\lambda_{q'} \circ f_1 = T(f_1) \circ \lambda_q$:
    $$
      \begin{tikzcd}
        E  \ar[r, "f_1"] \ar[d, "\lambda_q"'] & E' \ar[d, "\lambda_{q'}"] \\
        TE \ar[r, "T(f_1)"'] & TE'
      \end{tikzcd}
    $$
\end{enumerate}
\end{defn}

\begin{prop}[{\cite[Proposition~2.16]{CC18}}]\label{prop:linear-bundle-additive}
A linear bundle morphism is automatically \textbf{additive}.
That is, it preserves the addition and the zero-section of the differential bundle $\mathsf{q}$.
\end{prop}



 
\begin{thebibliography}{99}
\bibitem[AB]{AB} L. Aintablian, C. Blohmann, Differentiable groupoid objects and their
  abstract Lie algebroids, Appl. Categ. Structures 33, no. 5, Paper No. 33, 97 pp., 2025
\bibitem[Bloh23]{B23} C. Blohmann, Elastic diffeological spaces, Recent Advances in Diffeologies and Their Applications, Contemp. Math. 794, AMS, 2024, pp. 49–86.
\bibitem[Bloh24]{B24} C. Blohmann, Lagrangian Field Theory, Diffeology, Variational Cohomology,
  Multisymplectic Geometry, 2024. {\tt https://people.mpim-bonn.mpg.de/blohmann/LFT/}
\bibitem[CC14]{CC14} J.R.B. Cockett and G.S.H. Cruttwell. Differential structure, tangent
  structure, and SDG. Appl. Categ. Structures, {\bf 22} (2014), 331--417.
\bibitem[CC15]{CC15} J.R.B. Cockett and G.S.H. Cruttwell. The Jacobi identity for tangent
  categories. Cah. Topol. G\'eom. Diff\'er. Cat\'eg. {\bf 56} (2015), 301--316.
\bibitem[CC17]{CC17} J.R.B. Cockett and G.S.H. Cruttwell, Connections in tangent categories,
  Theory and Applications of Categories, {\bf 32} (2017), 835--888.
\bibitem[CC18]{CC18} J.R.B. Cockett and G.S.H. Cruttwell, Differential bundles and fibrations
  for tangent categories, Cah. Topol. G\'eom. Diff\'er. Cat\'eg. 59.1 (2018), pp.\ 10--92.
\bibitem[CCL]{CCL} J.R.B. Cockett, G.S.H. Cruttwell, J.S.P. Lemay, Differential equations in
  a tangent category I: Complete vector fields, flows, and exponentials,
  Applied Categorical Structures 29 (5), 773--825, 2021.
\bibitem[GW]{GW} N. Goldammer, K. Welker, Towards optimization techniques on diffeological spaces by generalizing Riemannian concepts, Appl. Math. Optim. 93, No. 1, Paper No. 5, 31 p. (2026). 
\bibitem[PIZ12]{PIZ12} P. Iglesias-Zemmour, Diffeology, Mathematical Surveys and Monographs,
  185, AMS, Providence, 2012.
\bibitem[PIZ25]{PIZ25} P. Iglesias-Zemmour, Lectures on diffeology,
  Beijing World Publishing Corporation, 2025.
\bibitem[Kihara]{Kihara} H. Kihara, Smooth homotopy of infinite-dimensional $C^\infty$-manifolds,
  Memoirs of the American Mathematical Society 1436. Providence, RI: American Mathematical Society (AMS), 2023.
\bibitem[KMS]{KMS} I. Kol\'a\v{r}, P. W. Michor, J. Slov\'ak, Natural Operations in Differential Geometry, Springer-Verlag Berlin, 1993.
\bibitem[KSS]{KSS} K. Kuribayashi, K. Sakai, Y. Shiobara, Towards Riemannian diffeology,
  Proc. Roy. Soc. Edinburgh Sect. A, published online 2025, pp.\ 1--30.
  \texttt{doi:10.1017/prm.2025.10114}.
\bibitem[LW]{LW} L.B.B. Lucyshyn-Wright, On the geometric notion of connection and its
  expression in tangent categories. Theory and Applications of Categories 33 (2018), 832--866.
\bibitem[Miya]{M} D. Miyamoto, Lie algebras of quotient groups, preprint, 2025.
  {\tt https://arxiv.org/abs/2502.10260v2}.
\bibitem[Ros]{R} J. Rosick\'y. Abstract tangent functors. Diagrammes, 12: JR1--JR11, 1984.
\bibitem[Taho]{Taho} M. Taho, Tangent spaces of diffeological spaces and their variants,
Topology and its Applications, Volume 381, 2026, {\tt https://doi.org/10.1016/j.topol.2026.109741}.
\bibitem[Vin]{V} M. Vincent, Diffeological differential geometry, MSc. Copenhagen, DK: University of Copenhagen, 2008.  
\end{thebibliography}
 
\end{document}
 
